% Options for packages loaded elsewhere
% Options for packages loaded elsewhere
\PassOptionsToPackage{unicode}{hyperref}
\PassOptionsToPackage{hyphens}{url}
\PassOptionsToPackage{dvipsnames,svgnames,x11names}{xcolor}
%
\documentclass[
  letterpaper,
  DIV=11,
  numbers=noendperiod]{scrartcl}
\usepackage{xcolor}
\usepackage{amsmath,amssymb}
\setcounter{secnumdepth}{-\maxdimen} % remove section numbering
\usepackage{iftex}
\ifPDFTeX
  \usepackage[T1]{fontenc}
  \usepackage[utf8]{inputenc}
  \usepackage{textcomp} % provide euro and other symbols
\else % if luatex or xetex
  \usepackage{unicode-math} % this also loads fontspec
  \defaultfontfeatures{Scale=MatchLowercase}
  \defaultfontfeatures[\rmfamily]{Ligatures=TeX,Scale=1}
\fi
\usepackage{lmodern}
\ifPDFTeX\else
  % xetex/luatex font selection
\fi
% Use upquote if available, for straight quotes in verbatim environments
\IfFileExists{upquote.sty}{\usepackage{upquote}}{}
\IfFileExists{microtype.sty}{% use microtype if available
  \usepackage[]{microtype}
  \UseMicrotypeSet[protrusion]{basicmath} % disable protrusion for tt fonts
}{}
\makeatletter
\@ifundefined{KOMAClassName}{% if non-KOMA class
  \IfFileExists{parskip.sty}{%
    \usepackage{parskip}
  }{% else
    \setlength{\parindent}{0pt}
    \setlength{\parskip}{6pt plus 2pt minus 1pt}}
}{% if KOMA class
  \KOMAoptions{parskip=half}}
\makeatother
% Make \paragraph and \subparagraph free-standing
\makeatletter
\ifx\paragraph\undefined\else
  \let\oldparagraph\paragraph
  \renewcommand{\paragraph}{
    \@ifstar
      \xxxParagraphStar
      \xxxParagraphNoStar
  }
  \newcommand{\xxxParagraphStar}[1]{\oldparagraph*{#1}\mbox{}}
  \newcommand{\xxxParagraphNoStar}[1]{\oldparagraph{#1}\mbox{}}
\fi
\ifx\subparagraph\undefined\else
  \let\oldsubparagraph\subparagraph
  \renewcommand{\subparagraph}{
    \@ifstar
      \xxxSubParagraphStar
      \xxxSubParagraphNoStar
  }
  \newcommand{\xxxSubParagraphStar}[1]{\oldsubparagraph*{#1}\mbox{}}
  \newcommand{\xxxSubParagraphNoStar}[1]{\oldsubparagraph{#1}\mbox{}}
\fi
\makeatother

\usepackage{color}
\usepackage{fancyvrb}
\newcommand{\VerbBar}{|}
\newcommand{\VERB}{\Verb[commandchars=\\\{\}]}
\DefineVerbatimEnvironment{Highlighting}{Verbatim}{commandchars=\\\{\}}
% Add ',fontsize=\small' for more characters per line
\usepackage{framed}
\definecolor{shadecolor}{RGB}{241,243,245}
\newenvironment{Shaded}{\begin{snugshade}}{\end{snugshade}}
\newcommand{\AlertTok}[1]{\textcolor[rgb]{0.68,0.00,0.00}{#1}}
\newcommand{\AnnotationTok}[1]{\textcolor[rgb]{0.37,0.37,0.37}{#1}}
\newcommand{\AttributeTok}[1]{\textcolor[rgb]{0.40,0.45,0.13}{#1}}
\newcommand{\BaseNTok}[1]{\textcolor[rgb]{0.68,0.00,0.00}{#1}}
\newcommand{\BuiltInTok}[1]{\textcolor[rgb]{0.00,0.23,0.31}{#1}}
\newcommand{\CharTok}[1]{\textcolor[rgb]{0.13,0.47,0.30}{#1}}
\newcommand{\CommentTok}[1]{\textcolor[rgb]{0.37,0.37,0.37}{#1}}
\newcommand{\CommentVarTok}[1]{\textcolor[rgb]{0.37,0.37,0.37}{\textit{#1}}}
\newcommand{\ConstantTok}[1]{\textcolor[rgb]{0.56,0.35,0.01}{#1}}
\newcommand{\ControlFlowTok}[1]{\textcolor[rgb]{0.00,0.23,0.31}{\textbf{#1}}}
\newcommand{\DataTypeTok}[1]{\textcolor[rgb]{0.68,0.00,0.00}{#1}}
\newcommand{\DecValTok}[1]{\textcolor[rgb]{0.68,0.00,0.00}{#1}}
\newcommand{\DocumentationTok}[1]{\textcolor[rgb]{0.37,0.37,0.37}{\textit{#1}}}
\newcommand{\ErrorTok}[1]{\textcolor[rgb]{0.68,0.00,0.00}{#1}}
\newcommand{\ExtensionTok}[1]{\textcolor[rgb]{0.00,0.23,0.31}{#1}}
\newcommand{\FloatTok}[1]{\textcolor[rgb]{0.68,0.00,0.00}{#1}}
\newcommand{\FunctionTok}[1]{\textcolor[rgb]{0.28,0.35,0.67}{#1}}
\newcommand{\ImportTok}[1]{\textcolor[rgb]{0.00,0.46,0.62}{#1}}
\newcommand{\InformationTok}[1]{\textcolor[rgb]{0.37,0.37,0.37}{#1}}
\newcommand{\KeywordTok}[1]{\textcolor[rgb]{0.00,0.23,0.31}{\textbf{#1}}}
\newcommand{\NormalTok}[1]{\textcolor[rgb]{0.00,0.23,0.31}{#1}}
\newcommand{\OperatorTok}[1]{\textcolor[rgb]{0.37,0.37,0.37}{#1}}
\newcommand{\OtherTok}[1]{\textcolor[rgb]{0.00,0.23,0.31}{#1}}
\newcommand{\PreprocessorTok}[1]{\textcolor[rgb]{0.68,0.00,0.00}{#1}}
\newcommand{\RegionMarkerTok}[1]{\textcolor[rgb]{0.00,0.23,0.31}{#1}}
\newcommand{\SpecialCharTok}[1]{\textcolor[rgb]{0.37,0.37,0.37}{#1}}
\newcommand{\SpecialStringTok}[1]{\textcolor[rgb]{0.13,0.47,0.30}{#1}}
\newcommand{\StringTok}[1]{\textcolor[rgb]{0.13,0.47,0.30}{#1}}
\newcommand{\VariableTok}[1]{\textcolor[rgb]{0.07,0.07,0.07}{#1}}
\newcommand{\VerbatimStringTok}[1]{\textcolor[rgb]{0.13,0.47,0.30}{#1}}
\newcommand{\WarningTok}[1]{\textcolor[rgb]{0.37,0.37,0.37}{\textit{#1}}}

\usepackage{longtable,booktabs,array}
\newcounter{none} % for unnumbered tables
\usepackage{calc} % for calculating minipage widths
% Correct order of tables after \paragraph or \subparagraph
\usepackage{etoolbox}
\makeatletter
\patchcmd\longtable{\par}{\if@noskipsec\mbox{}\fi\par}{}{}
\makeatother
% Allow footnotes in longtable head/foot
\IfFileExists{footnotehyper.sty}{\usepackage{footnotehyper}}{\usepackage{footnote}}
\makesavenoteenv{longtable}
\usepackage{graphicx}
\makeatletter
\newsavebox\pandoc@box
\newcommand*\pandocbounded[1]{% scales image to fit in text height/width
  \sbox\pandoc@box{#1}%
  \Gscale@div\@tempa{\textheight}{\dimexpr\ht\pandoc@box+\dp\pandoc@box\relax}%
  \Gscale@div\@tempb{\linewidth}{\wd\pandoc@box}%
  \ifdim\@tempb\p@<\@tempa\p@\let\@tempa\@tempb\fi% select the smaller of both
  \ifdim\@tempa\p@<\p@\scalebox{\@tempa}{\usebox\pandoc@box}%
  \else\usebox{\pandoc@box}%
  \fi%
}
% Set default figure placement to htbp
\def\fps@figure{htbp}
\makeatother





\setlength{\emergencystretch}{3em} % prevent overfull lines

\providecommand{\tightlist}{%
  \setlength{\itemsep}{0pt}\setlength{\parskip}{0pt}}



 


\KOMAoption{captions}{tableheading}
\makeatletter
\@ifpackageloaded{caption}{}{\usepackage{caption}}
\AtBeginDocument{%
\ifdefined\contentsname
  \renewcommand*\contentsname{Table of contents}
\else
  \newcommand\contentsname{Table of contents}
\fi
\ifdefined\listfigurename
  \renewcommand*\listfigurename{List of Figures}
\else
  \newcommand\listfigurename{List of Figures}
\fi
\ifdefined\listtablename
  \renewcommand*\listtablename{List of Tables}
\else
  \newcommand\listtablename{List of Tables}
\fi
\ifdefined\figurename
  \renewcommand*\figurename{Figure}
\else
  \newcommand\figurename{Figure}
\fi
\ifdefined\tablename
  \renewcommand*\tablename{Table}
\else
  \newcommand\tablename{Table}
\fi
}
\@ifpackageloaded{float}{}{\usepackage{float}}
\floatstyle{ruled}
\@ifundefined{c@chapter}{\newfloat{codelisting}{h}{lop}}{\newfloat{codelisting}{h}{lop}[chapter]}
\floatname{codelisting}{Listing}
\newcommand*\listoflistings{\listof{codelisting}{List of Listings}}
\makeatother
\makeatletter
\makeatother
\makeatletter
\@ifpackageloaded{caption}{}{\usepackage{caption}}
\@ifpackageloaded{subcaption}{}{\usepackage{subcaption}}
\makeatother
\usepackage{bookmark}
\IfFileExists{xurl.sty}{\usepackage{xurl}}{} % add URL line breaks if available
\urlstyle{same}
\hypersetup{
  colorlinks=true,
  linkcolor={blue},
  filecolor={Maroon},
  citecolor={Blue},
  urlcolor={Blue},
  pdfcreator={LaTeX via pandoc}}


\author{}
\date{}
\begin{document}


\section{Regressão Logística
Binária}\label{regressuxe3o-loguxedstica-binuxe1ria}

\emph{Tutorial Completo --- Teoria, Formalismo e Implementação em R e
Python}

O tutorial nasceu de notas de estudos feitas manuscritas. As leituras
que embasaram este estudo são de livros de estatística e ciência de
dados que serão referenciadas nesta apresentação.

O manuscrito que deu origem a esta apresentação foi
\href{analise_de_regressao_logistica_versao_1.pdf}{Manuscrito Original}

Caso esta apresentação não este renderizando de forma adequada use o pdf
gerado ou o html para realizar a leitura!

\textbf{Autor:} Paulo Pimenta

\begin{center}\rule{0.5\linewidth}{0.5pt}\end{center}

\begin{quote}
\textbf{Sobre este tutorial}\\
Cobertura completa da regressão logística binária: da motivação
matemática até a implementação prática em \textbf{R} e \textbf{Python}.
O código aceita qualquer arquivo \texttt{.csv} --- basta ajustar o bloco
de configuração indicado em cada seção.\\
Baseado em notas de aula revisadas, com formalismo estatístico e
conteúdos complementares ausentes no manuscrito original.
\end{quote}

\begin{center}\rule{0.5\linewidth}{0.5pt}\end{center}

\section{1. Introdução e
Motivação}\label{introduuxe7uxe3o-e-motivauxe7uxe3o}

\subsection{1.1 Por que não usar regressão
linear?}\label{por-que-nuxe3o-usar-regressuxe3o-linear}

Quando a variável resposta \(y\) é \textbf{binária} --- assume apenas os
valores \(0\) (fracasso, ausência, não-evento) e \(1\) (sucesso,
presença, evento) --- a regressão linear ordinária apresenta problemas
fundamentais:

{\def\LTcaptype{none} % do not increment counter
\begin{longtable}[]{@{}
  >{\raggedright\arraybackslash}p{(\linewidth - 2\tabcolsep) * \real{0.5000}}
  >{\raggedright\arraybackslash}p{(\linewidth - 2\tabcolsep) * \real{0.5000}}@{}}
\toprule\noalign{}
\begin{minipage}[b]{\linewidth}\raggedright
Problema
\end{minipage} & \begin{minipage}[b]{\linewidth}\raggedright
Descrição
\end{minipage} \\
\midrule\noalign{}
\endhead
\bottomrule\noalign{}
\endlastfoot
\textbf{Previsões fora de {[}0,1{]}} &
\(\hat{y} = \mathbf{x}^\top\pmb{\beta}\) pode assumir qualquer valor
real, resultando em ``probabilidades'' negativas ou maiores que 1 \\
\textbf{Heterocedasticidade estrutural} & A variância de
\(y \sim \text{Bernoulli}(p)\) é \(p(1-p)\), que varia com \(p\),
violando a homocedasticidade \\
\textbf{Distribuição dos resíduos} & Os resíduos não seguem distribuição
normal, invalidando os testes da regressão linear \\
\end{longtable}
}

A solução é modelar diretamente a \textbf{probabilidade condicional}
\(P(Y=1 \mid \mathbf{x})\) por meio de uma função que mapeie
\((-\infty, +\infty) \to (0, 1)\). A função escolhida é a
\textbf{sigmoide logística}.

\subsection{1.2 Aplicações típicas}\label{aplicauxe7uxf5es-tuxedpicas}

\begin{itemize}
\tightlist
\item
  \textbf{Medicina:} diagnóstico (doente/saudável), presença de doença
  (sim/não)\\
\item
  \textbf{Finanças:} inadimplência (sim/não), fraude (sim/não)\\
\item
  \textbf{Marketing:} conversão de cliente (compra/não compra),
  \emph{churn} (cancela/permanece)\\
\item
  \textbf{NLP:} detecção de spam (spam/não-spam), sentimento
  (positivo/negativo)\\
\item
  \textbf{Biologia:} sobrevivência de espécie (sobrevive/extingue)
\end{itemize}

\subsection{1.3 Fluxo da análise}\label{fluxo-da-anuxe1lise}

O pipeline completo de uma análise de regressão logística binária:

\begin{verbatim}
Dados .csv
    │
    ▼
① Exploração (dispersão, boxplots, balanço de classes)
    │
    ▼
② Estimação do modelo (Máxima Verossimilhança)
    │
    ▼
③ Avaliação da acurácia (pseudo-R², acurácia, AUC-ROC)
    │
    ▼
④ Testes de hipóteses (TRV global + Wald individual)
    │
    ▼
⑤ Previsão para novas observações
\end{verbatim}

\begin{center}\rule{0.5\linewidth}{0.5pt}\end{center}

\section{2. O Modelo Logístico}\label{o-modelo-loguxedstico}

\subsection{2.1 A função sigmoide}\label{a-funuxe7uxe3o-sigmoide}

A equação central do modelo é:

\[\hat{y} = P(Y = 1 \mid \mathbf{x}) = \frac{1}{1 + e^{-z}}, \quad z = a_1 x_1 + a_2 x_2 + \cdots + a_p x_p + b\]

onde:

{\def\LTcaptype{none} % do not increment counter
\begin{longtable}[]{@{}
  >{\centering\arraybackslash}p{(\linewidth - 4\tabcolsep) * \real{0.3846}}
  >{\raggedright\arraybackslash}p{(\linewidth - 4\tabcolsep) * \real{0.3077}}
  >{\raggedright\arraybackslash}p{(\linewidth - 4\tabcolsep) * \real{0.3077}}@{}}
\toprule\noalign{}
\begin{minipage}[b]{\linewidth}\centering
Símbolo
\end{minipage} & \begin{minipage}[b]{\linewidth}\raggedright
Nome
\end{minipage} & \begin{minipage}[b]{\linewidth}\raggedright
Descrição
\end{minipage} \\
\midrule\noalign{}
\endhead
\bottomrule\noalign{}
\endlastfoot
\(x_1, \ldots, x_p\) & Variáveis preditoras & Características
observadas \\
\(a_1, \ldots, a_p\) & Coeficientes de regressão & Parâmetros estimados
pelo modelo \\
\(b\) & Intercepto & Parâmetro constante (viés) \\
\(z\) & Log-odds (logit) & Combinação linear das preditoras \\
\(\hat{y}\) & Probabilidade predita &
\(P(Y=1 \mid \mathbf{x}) \in (0,1)\) \\
\end{longtable}
}

\textbf{Propriedade fundamental:} independentemente do valor de \(z\),
\(0 < \hat{y} < 1\) sempre. O modelo nunca produz uma probabilidade
inválida.

\begin{Shaded}
\begin{Highlighting}[]
\FunctionTok{library}\NormalTok{(ggplot2)}
\FunctionTok{library}\NormalTok{(dplyr)}
\FunctionTok{library}\NormalTok{(pROC)}
\FunctionTok{library}\NormalTok{(caret)}
\FunctionTok{library}\NormalTok{(gridExtra)}

\NormalTok{z\_vals }\OtherTok{\textless{}{-}} \FunctionTok{seq}\NormalTok{(}\SpecialCharTok{{-}}\DecValTok{10}\NormalTok{, }\DecValTok{10}\NormalTok{, }\AttributeTok{length.out =} \DecValTok{600}\NormalTok{)}
\NormalTok{sig\_df }\OtherTok{\textless{}{-}} \FunctionTok{data.frame}\NormalTok{(}\AttributeTok{z =}\NormalTok{ z\_vals, }\AttributeTok{y =} \DecValTok{1} \SpecialCharTok{/}\NormalTok{ (}\DecValTok{1} \SpecialCharTok{+} \FunctionTok{exp}\NormalTok{(}\SpecialCharTok{{-}}\NormalTok{z\_vals)))}

\FunctionTok{ggplot}\NormalTok{(sig\_df, }\FunctionTok{aes}\NormalTok{(}\AttributeTok{x =}\NormalTok{ z, }\AttributeTok{y =}\NormalTok{ y)) }\SpecialCharTok{+}
  \FunctionTok{geom\_ribbon}\NormalTok{(}\AttributeTok{data =} \FunctionTok{filter}\NormalTok{(sig\_df, y }\SpecialCharTok{\textgreater{}=} \FloatTok{0.5}\NormalTok{),}
              \FunctionTok{aes}\NormalTok{(}\AttributeTok{ymin =} \FloatTok{0.5}\NormalTok{, }\AttributeTok{ymax =}\NormalTok{ y), }\AttributeTok{fill =} \StringTok{"\#27AE60"}\NormalTok{, }\AttributeTok{alpha =} \FloatTok{0.18}\NormalTok{) }\SpecialCharTok{+}
  \FunctionTok{geom\_ribbon}\NormalTok{(}\AttributeTok{data =} \FunctionTok{filter}\NormalTok{(sig\_df, y }\SpecialCharTok{\textless{}} \FloatTok{0.5}\NormalTok{),}
              \FunctionTok{aes}\NormalTok{(}\AttributeTok{ymin =}\NormalTok{ y, }\AttributeTok{ymax =} \FloatTok{0.5}\NormalTok{), }\AttributeTok{fill =} \StringTok{"\#C0392B"}\NormalTok{, }\AttributeTok{alpha =} \FloatTok{0.18}\NormalTok{) }\SpecialCharTok{+}
  \FunctionTok{geom\_line}\NormalTok{(}\AttributeTok{color =} \StringTok{"\#2C3E50"}\NormalTok{, }\AttributeTok{size =} \FloatTok{1.4}\NormalTok{) }\SpecialCharTok{+}
  \FunctionTok{geom\_hline}\NormalTok{(}\AttributeTok{yintercept =} \FunctionTok{c}\NormalTok{(}\DecValTok{0}\NormalTok{, }\FloatTok{0.5}\NormalTok{, }\DecValTok{1}\NormalTok{), }\AttributeTok{linetype =} \StringTok{"dashed"}\NormalTok{,}
             \AttributeTok{color =} \FunctionTok{c}\NormalTok{(}\StringTok{"gray70"}\NormalTok{,}\StringTok{"gray40"}\NormalTok{,}\StringTok{"gray70"}\NormalTok{), }\AttributeTok{size =} \FloatTok{0.6}\NormalTok{) }\SpecialCharTok{+}
  \FunctionTok{geom\_vline}\NormalTok{(}\AttributeTok{xintercept =} \DecValTok{0}\NormalTok{, }\AttributeTok{linetype =} \StringTok{"dashed"}\NormalTok{, }\AttributeTok{color =} \StringTok{"gray40"}\NormalTok{, }\AttributeTok{size =} \FloatTok{0.6}\NormalTok{) }\SpecialCharTok{+}
  \FunctionTok{geom\_point}\NormalTok{(}\AttributeTok{data =} \FunctionTok{data.frame}\NormalTok{(}\AttributeTok{z =} \DecValTok{0}\NormalTok{, }\AttributeTok{y =} \FloatTok{0.5}\NormalTok{),}
             \FunctionTok{aes}\NormalTok{(}\AttributeTok{x =}\NormalTok{ z, }\AttributeTok{y =}\NormalTok{ y), }\AttributeTok{color =} \StringTok{"\#E74C3C"}\NormalTok{, }\AttributeTok{size =} \DecValTok{3}\NormalTok{) }\SpecialCharTok{+}
  \FunctionTok{annotate}\NormalTok{(}\StringTok{"text"}\NormalTok{, }\AttributeTok{x =} \FloatTok{6.5}\NormalTok{, }\AttributeTok{y =} \FloatTok{0.18}\NormalTok{,}
           \AttributeTok{label =} \StringTok{"hat(y) == frac(1, 1 + e\^{}\{{-}z\})"}\NormalTok{,}
           \AttributeTok{parse =} \ConstantTok{TRUE}\NormalTok{, }\AttributeTok{size =} \FloatTok{4.5}\NormalTok{, }\AttributeTok{color =} \StringTok{"\#2C3E50"}\NormalTok{) }\SpecialCharTok{+}
  \FunctionTok{annotate}\NormalTok{(}\StringTok{"text"}\NormalTok{, }\AttributeTok{x =} \DecValTok{4}\NormalTok{, }\AttributeTok{y =} \FloatTok{0.80}\NormalTok{, }\AttributeTok{label =} \StringTok{"Classe 1"}\NormalTok{, }\AttributeTok{color =} \StringTok{"\#27AE60"}\NormalTok{,}
           \AttributeTok{fontface =} \StringTok{"bold"}\NormalTok{, }\AttributeTok{size =} \DecValTok{4}\NormalTok{) }\SpecialCharTok{+}
  \FunctionTok{annotate}\NormalTok{(}\StringTok{"text"}\NormalTok{, }\AttributeTok{x =} \SpecialCharTok{{-}}\DecValTok{4}\NormalTok{, }\AttributeTok{y =} \FloatTok{0.18}\NormalTok{, }\AttributeTok{label =} \StringTok{"Classe 0"}\NormalTok{, }\AttributeTok{color =} \StringTok{"\#C0392B"}\NormalTok{,}
           \AttributeTok{fontface =} \StringTok{"bold"}\NormalTok{, }\AttributeTok{size =} \DecValTok{4}\NormalTok{) }\SpecialCharTok{+}
  \FunctionTok{scale\_y\_continuous}\NormalTok{(}\AttributeTok{breaks =} \FunctionTok{c}\NormalTok{(}\DecValTok{0}\NormalTok{, }\FloatTok{0.25}\NormalTok{, }\FloatTok{0.5}\NormalTok{, }\FloatTok{0.75}\NormalTok{, }\DecValTok{1}\NormalTok{),}
                     \AttributeTok{labels =} \FunctionTok{c}\NormalTok{(}\StringTok{"0"}\NormalTok{,}\StringTok{"0,25"}\NormalTok{,}\StringTok{"0,5"}\NormalTok{,}\StringTok{"0,75"}\NormalTok{,}\StringTok{"1"}\NormalTok{)) }\SpecialCharTok{+}
  \FunctionTok{scale\_x\_continuous}\NormalTok{(}\AttributeTok{breaks =} \FunctionTok{seq}\NormalTok{(}\SpecialCharTok{{-}}\DecValTok{10}\NormalTok{, }\DecValTok{10}\NormalTok{, }\DecValTok{2}\NormalTok{)) }\SpecialCharTok{+}
  \FunctionTok{labs}\NormalTok{(}\AttributeTok{x =} \StringTok{"z  (log{-}odds)"}\NormalTok{, }\AttributeTok{y =} \FunctionTok{expression}\NormalTok{(}\FunctionTok{hat}\NormalTok{(y) }\SpecialCharTok{==} \FunctionTok{P}\NormalTok{(Y }\SpecialCharTok{==} \DecValTok{1} \SpecialCharTok{\textasciitilde{}} \StringTok{"|"} \SpecialCharTok{\textasciitilde{}} \FunctionTok{bold}\NormalTok{(x))),}
       \AttributeTok{title =} \StringTok{"Função Logística (Sigmoide)"}\NormalTok{) }\SpecialCharTok{+}
  \FunctionTok{theme\_minimal}\NormalTok{(}\AttributeTok{base\_size =} \DecValTok{13}\NormalTok{) }\SpecialCharTok{+}
  \FunctionTok{theme}\NormalTok{(}\AttributeTok{panel.grid.minor =} \FunctionTok{element\_blank}\NormalTok{())}
\end{Highlighting}
\end{Shaded}

\subsection{2.2 A transformação logit e os
odds}\label{a-transformauxe7uxe3o-logit-e-os-odds}

Invertendo a função sigmoide:

\[\log_e\!\left(\frac{p}{1-p}\right) \; [\text{logit}(p)] = z = a_1 x_1 + \cdots + a_p x_p + b\]

O termo \(\dfrac{p}{1-p}\) é a \textbf{razão de chances} (\emph{odds}):
razão entre a probabilidade de ocorrência e de não-ocorrência do evento.
A transformação logit lineariza a relação entre as preditoras e os odds
do evento, que é a base para interpretar os coeficientes.

\subsubsection{Interpretação dos coeficientes via Odds
Ratio}\label{interpretauxe7uxe3o-dos-coeficientes-via-odds-ratio}

\[OR_i = e^{a_i}\]

{\def\LTcaptype{none} % do not increment counter
\begin{longtable}[]{@{}
  >{\raggedright\arraybackslash}p{(\linewidth - 2\tabcolsep) * \real{0.5000}}
  >{\raggedright\arraybackslash}p{(\linewidth - 2\tabcolsep) * \real{0.5000}}@{}}
\toprule\noalign{}
\begin{minipage}[b]{\linewidth}\raggedright
Situação
\end{minipage} & \begin{minipage}[b]{\linewidth}\raggedright
Interpretação
\end{minipage} \\
\midrule\noalign{}
\endhead
\bottomrule\noalign{}
\endlastfoot
\(a_i > 0\) \(\Rightarrow\) \(OR_i > 1\) & Aumento de 1 unidade em
\(x_i\) \textbf{multiplica} os odds por \(e^{a_i}\) \\
\(a_i < 0\) \(\Rightarrow\) \(OR_i < 1\) & Aumento de 1 unidade em
\(x_i\) \textbf{reduz} os odds por fator \(e^{a_i}\) \\
\(a_i = 0\) \(\Rightarrow\) \(OR_i = 1\) & \(x_i\) não tem efeito sobre
os odds do evento \\
\end{longtable}
}

\begin{quote}
\textbf{Exemplo numérico:} se \(a_1 = 2{,}44\), então
\(OR_1 = e^{2,44} \approx 11{,}47\). Um aumento de 1 unidade em \(x_1\)
eleva os odds do evento em \textbf{1047\%}, mantidas as demais variáveis
constantes.
\end{quote}

\begin{center}\rule{0.5\linewidth}{0.5pt}\end{center}

\section{3. Estimação por Máxima
Verossimilhança}\label{estimauxe7uxe3o-por-muxe1xima-verossimilhanuxe7a}

\subsection{3.1 O princípio}\label{o-princuxedpio}

A \textbf{Máxima Verossimilhança} (MV) encontra os valores de
\(\mathbf{a}\) e \(b\) que \textbf{maximizam a probabilidade de se
observar exatamente os dados coletados}, dada a estrutura do modelo.

\subsection{3.2 A função de
verossimilhança}\label{a-funuxe7uxe3o-de-verossimilhanuxe7a}

Para \(n\) observações independentes, com \(y_i \in \{0,1\}\), cada
observação segue uma distribuição de Bernoulli com parâmetro
\(\hat{y}_i\). A verossimilhança conjunta é:

\[\mathcal{L}(\mathbf{a}, b) = \prod_{i=1}^{n} \hat{y}_i^{\;y_i} \cdot (1 - \hat{y}_i)^{1-y_i}\]

Cada fator contribui com: - \(\hat{y}_i\) quando \(y_i = 1\)
(probabilidade prevista de sucesso)\\
- \((1 - \hat{y}_i)\) quando \(y_i = 0\) (probabilidade prevista de
fracasso)

\subsection{3.3 A log-verossimilhança}\label{a-log-verossimilhanuxe7a}

Para evitar \emph{underflow} numérico e transformar o produto em soma,
maximiza-se o logaritmo natural:

\[L(\mathbf{a}, b) = \sum_{i=1}^{n}\left[\, y_i \log_e(\hat{y}_i) + (1-y_i)\log_e(1-\hat{y}_i) \,\right]\]

\textbf{Propriedades:}

\begin{itemize}
\tightlist
\item
  \(L \leq 0\) sempre, pois \(\log_e(\hat{y}_i) \leq 0\) para
  \(\hat{y}_i \in (0,1)\)
\item
  \(L = 0\) somente quando o modelo classifica perfeitamente todos os
  pontos
\item
  Maximizar \(L\) equivale a minimizar a \textbf{entropia cruzada
  binária}, função de perda usada em redes neurais para classificação
  binária
\end{itemize}

\subsection{3.4 Exemplo: preferência por
café}\label{exemplo-preferuxeancia-por-cafuxe9}

Pesquisa com 10 pessoas, das quais 7 responderam ``Sim'' e 3 ``Não''.
Parâmetro a estimar: \(p\) (proporção que gosta de café).

\textbf{Verossimilhança} (desconsiderando combinatórias --- não afetam o
maximizador): \[\mathcal{L}(p) = p^7 \cdot (1-p)^3\]

\textbf{Log-verossimilhança:} \[L(p) = 7\log_e(p) + 3\log_e(1-p)\]

\textbf{Condição de 1ª ordem} (ponto crítico):
\[\frac{dL}{dp} = \frac{7}{p} - \frac{3}{1-p} = 0 \;\implies\; 7(1-p) = 3p \;\implies\; \hat{p} = \frac{7}{10} = 0{,}7\]

\textbf{Condição de 2ª ordem} (verificação de máximo):
\[\frac{d^2L}{dp^2} = -\frac{7}{p^2} - \frac{3}{(1-p)^2} < 0 \quad \forall\, p \in (0,1)\]

Como a segunda derivada é estritamente negativa, \(\hat{p} = 0{,}7\) é
um \textbf{máximo global} --- coincide, como esperado, com a proporção
amostral.

\begin{Shaded}
\begin{Highlighting}[]
\NormalTok{p\_seq  }\OtherTok{\textless{}{-}} \FunctionTok{seq}\NormalTok{(}\FloatTok{0.01}\NormalTok{, }\FloatTok{0.99}\NormalTok{, }\AttributeTok{length.out =} \DecValTok{500}\NormalTok{)}
\NormalTok{L\_cafe }\OtherTok{\textless{}{-}} \DecValTok{7} \SpecialCharTok{*} \FunctionTok{log}\NormalTok{(p\_seq) }\SpecialCharTok{+} \DecValTok{3} \SpecialCharTok{*} \FunctionTok{log}\NormalTok{(}\DecValTok{1} \SpecialCharTok{{-}}\NormalTok{ p\_seq)}
\NormalTok{p\_hat  }\OtherTok{\textless{}{-}} \DecValTok{7} \SpecialCharTok{/} \DecValTok{10}
\NormalTok{L\_hat  }\OtherTok{\textless{}{-}} \DecValTok{7} \SpecialCharTok{*} \FunctionTok{log}\NormalTok{(p\_hat) }\SpecialCharTok{+} \DecValTok{3} \SpecialCharTok{*} \FunctionTok{log}\NormalTok{(}\DecValTok{1} \SpecialCharTok{{-}}\NormalTok{ p\_hat)}

\NormalTok{mv\_df  }\OtherTok{\textless{}{-}} \FunctionTok{data.frame}\NormalTok{(}\AttributeTok{p =}\NormalTok{ p\_seq, }\AttributeTok{L =}\NormalTok{ L\_cafe)}

\FunctionTok{ggplot}\NormalTok{(mv\_df, }\FunctionTok{aes}\NormalTok{(}\AttributeTok{x =}\NormalTok{ p, }\AttributeTok{y =}\NormalTok{ L)) }\SpecialCharTok{+}
  \FunctionTok{geom\_line}\NormalTok{(}\AttributeTok{color =} \StringTok{"\#8E44AD"}\NormalTok{, }\AttributeTok{size =} \FloatTok{1.3}\NormalTok{) }\SpecialCharTok{+}
  \FunctionTok{geom\_vline}\NormalTok{(}\AttributeTok{xintercept =}\NormalTok{ p\_hat, }\AttributeTok{linetype =} \StringTok{"dashed"}\NormalTok{,}
             \AttributeTok{color =} \StringTok{"\#C0392B"}\NormalTok{, }\AttributeTok{size =} \FloatTok{0.8}\NormalTok{) }\SpecialCharTok{+}
  \FunctionTok{geom\_hline}\NormalTok{(}\AttributeTok{yintercept =}\NormalTok{ L\_hat, }\AttributeTok{linetype =} \StringTok{"dotted"}\NormalTok{,}
             \AttributeTok{color =} \StringTok{"gray50"}\NormalTok{, }\AttributeTok{size =} \FloatTok{0.7}\NormalTok{) }\SpecialCharTok{+}
  \FunctionTok{geom\_point}\NormalTok{(}\FunctionTok{aes}\NormalTok{(}\AttributeTok{x =}\NormalTok{ p\_hat, }\AttributeTok{y =}\NormalTok{ L\_hat),}
             \AttributeTok{color =} \StringTok{"\#C0392B"}\NormalTok{, }\AttributeTok{size =} \DecValTok{4}\NormalTok{, }\AttributeTok{shape =} \DecValTok{21}\NormalTok{,}
             \AttributeTok{fill =} \StringTok{"\#E74C3C"}\NormalTok{) }\SpecialCharTok{+}
  \FunctionTok{annotate}\NormalTok{(}\StringTok{"text"}\NormalTok{, }\AttributeTok{x =} \FloatTok{0.76}\NormalTok{, }\AttributeTok{y =}\NormalTok{ L\_hat }\SpecialCharTok{{-}} \FloatTok{0.3}\NormalTok{,}
           \AttributeTok{label =} \FunctionTok{expression}\NormalTok{(}\FunctionTok{hat}\NormalTok{(p) }\SpecialCharTok{==} \FloatTok{0.7}\NormalTok{),}
           \AttributeTok{color =} \StringTok{"\#C0392B"}\NormalTok{, }\AttributeTok{size =} \DecValTok{5}\NormalTok{, }\AttributeTok{fontface =} \StringTok{"bold"}\NormalTok{) }\SpecialCharTok{+}
  \FunctionTok{annotate}\NormalTok{(}\StringTok{"text"}\NormalTok{, }\AttributeTok{x =} \FloatTok{0.20}\NormalTok{, }\AttributeTok{y =}\NormalTok{ L\_hat,}
           \AttributeTok{label =} \FunctionTok{paste0}\NormalTok{(}\StringTok{"L* = "}\NormalTok{, }\FunctionTok{round}\NormalTok{(L\_hat, }\DecValTok{3}\NormalTok{)),}
           \AttributeTok{color =} \StringTok{"gray40"}\NormalTok{, }\AttributeTok{size =} \DecValTok{4}\NormalTok{, }\AttributeTok{vjust =} \SpecialCharTok{{-}}\FloatTok{0.6}\NormalTok{) }\SpecialCharTok{+}
  \FunctionTok{labs}\NormalTok{(}\AttributeTok{x =} \StringTok{"p"}\NormalTok{, }\AttributeTok{y =} \StringTok{"L(p)"}\NormalTok{,}
       \AttributeTok{title =} \StringTok{"Log{-}Verossimilhança — Exemplo do Café"}\NormalTok{,}
       \AttributeTok{subtitle =} \StringTok{"Máximo em p̂ = 7/10 = 0,7"}\NormalTok{) }\SpecialCharTok{+}
  \FunctionTok{theme\_minimal}\NormalTok{(}\AttributeTok{base\_size =} \DecValTok{13}\NormalTok{) }\SpecialCharTok{+}
  \FunctionTok{theme}\NormalTok{(}\AttributeTok{panel.grid.minor =} \FunctionTok{element\_blank}\NormalTok{())}
\end{Highlighting}
\end{Shaded}

\subsection{3.5 Log-verossimilhança com preditoras (modelo
geral)}\label{log-verossimilhanuxe7a-com-preditoras-modelo-geral}

Substituindo \(\hat{y}_i = \sigma(z_i)\):

\[L(\mathbf{a}, b) = \sum_{i=1}^{n} \left[\, y_i \log_e\!\left(\frac{1}{1+e^{-z_i}}\right) + (1-y_i)\log_e\!\left(\frac{e^{-z_i}}{1+e^{-z_i}}\right) \,\right]\]

\textbf{Não existe solução analítica fechada.} Os coeficientes são
estimados por algoritmos de otimização iterativa:

{\def\LTcaptype{none} % do not increment counter
\begin{longtable}[]{@{}
  >{\raggedright\arraybackslash}p{(\linewidth - 4\tabcolsep) * \real{0.3333}}
  >{\raggedright\arraybackslash}p{(\linewidth - 4\tabcolsep) * \real{0.3333}}
  >{\raggedright\arraybackslash}p{(\linewidth - 4\tabcolsep) * \real{0.3333}}@{}}
\toprule\noalign{}
\begin{minipage}[b]{\linewidth}\raggedright
Algoritmo
\end{minipage} & \begin{minipage}[b]{\linewidth}\raggedright
Descrição
\end{minipage} & \begin{minipage}[b]{\linewidth}\raggedright
Implementação
\end{minipage} \\
\midrule\noalign{}
\endhead
\bottomrule\noalign{}
\endlastfoot
Newton-Raphson (IRLS) & Usa gradiente e hessiana exata; converge
rapidamente & \texttt{glm()} no R \\
BFGS / L-BFGS & Aproxima a hessiana; eficiente para muitas variáveis &
\texttt{statsmodels}, \texttt{sklearn} \\
Gradiente Descendente & Atualização por mini-batch; escalável a Big Data
& Redes neurais, TensorFlow \\
\end{longtable}
}

\begin{center}\rule{0.5\linewidth}{0.5pt}\end{center}

\section{4. Avaliação do Modelo}\label{avaliauxe7uxe3o-do-modelo}

\subsection{4.1 Pseudo-R² de McFadden}\label{pseudo-ruxb2-de-mcfadden}

Na regressão logística não existe um \(R^2\) com interpretação
geométrica direta. O \textbf{pseudo-\(R^2\) de McFadden} (1974) é a
medida mais usada:

\[R^2_{\text{McFadden}} = 1 -\frac{L^*}{L_0} = 1 - \frac{L^*}{N_1\log_e N_1 + N_0\log_e N_0 - (N_1+N_0)\log_e(N_1+N_0)}\]

onde:

{\def\LTcaptype{none} % do not increment counter
\begin{longtable}[]{@{}cl@{}}
\toprule\noalign{}
Símbolo & Definição \\
\midrule\noalign{}
\endhead
\bottomrule\noalign{}
\endlastfoot
\(L^*\) & Log-verossimilhança máxima do modelo ajustado (\(\leq 0\)) \\
\(L_0\) & Log-verossimilhança do modelo nulo (só intercepto) \\
\(N_1\) & Número de sucessos (\(y = 1\)) \\
\(N_0\) & Número de fracassos (\(y = 0\)) \\
\end{longtable}
}

O denominador \(L_0\) equivale à log-verossimilhança de um modelo que
prevê \(\hat{p} = N_1/N\) para todos --- o ``pior'' modelo útil.

\textbf{Tabela de interpretação:}

{\def\LTcaptype{none} % do not increment counter
\begin{longtable}[]{@{}cl@{}}
\toprule\noalign{}
Valor de \(R^2_{\text{McFadden}}\) & Avaliação \\
\midrule\noalign{}
\endhead
\bottomrule\noalign{}
\endlastfoot
\(< 0{,}20\) & Ajuste fraco \\
\(0{,}20 - 0{,}40\) & Bom ajuste \\
\(> 0{,}40\) & Ajuste muito bom \\
\end{longtable}
}

\begin{quote}
Na regressão logística, valores relativamente baixos de pseudo-\(R^2\)
são esperados e não indicam necessariamente um modelo ruim. Um
\(R^2_{\text{McFadden}} \approx 0{,}4\) em regressão logística é
comparável a um \(R^2 \approx 0{,}9\) em regressão linear.
\end{quote}

\subsection{4.1 R² de McFadden -- Fórmula
detalhada}\label{ruxb2-de-mcfadden-fuxf3rmula-detalhada}

\subsubsection{4.1.1 Fórmula correta
(conceitual)}\label{fuxf3rmula-correta-conceitual}

\[
R^2_{\text{McFadden}} = 1 - \frac{\ln L^*}{\ln L_0}
\]

\begin{itemize}
\tightlist
\item
  \(\ln L^*\) = log-verossimilhança do \textbf{modelo completo} (com
  preditores).\\
\item
  \(\ln L_0\) = log-verossimilhança do \textbf{modelo nulo} (apenas
  intercepto).
\end{itemize}

Ambos são números \textbf{negativos} (ou zero).\\
Como o modelo completo é pelo menos tão bom quanto o nulo,
\[\ln L^* \ge \ln L_0\] (menos negativo).\\
Portanto \[
0 \le \frac{\ln L^*}{\ln L_0} \le 1
\] e o \(R^2\) fica entre 0 e 1.

\begin{quote}
O erro comum é omitir o termo \(1 - \cdots\), escrevendo apenas
\[-\frac{L^*}{L_0}\], o que gera valor e sinal incorretos.
\end{quote}

\begin{center}\rule{0.5\linewidth}{0.5pt}\end{center}

\subsubsection{\texorpdfstring{4.1.2 Fórmula detalhada do modelo nulo
(\(L_0\))}{4.1.2 Fórmula detalhada do modelo nulo (L\_0)}}\label{fuxf3rmula-detalhada-do-modelo-nulo-l_0}

Em regressão logística binária (\(Y \in \{0,1\}\)), o modelo nulo prevê
a mesma probabilidade para todos:\\
\(\hat{p} = \frac{N_1}{N}\), onde:

\begin{itemize}
\tightlist
\item
  \(N_1\) = observações com \(Y=1\)\\
\item
  \(N_0\) = observações com \(Y=0\)\\
\item
  \(N = N_1 + N_0\)
\end{itemize}

A log-verossimilhança do modelo nulo é:

\[
\ln L_0 = N_1 \ln\!\left(\frac{N_1}{N}\right) + N_0 \ln\!\left(\frac{N_0}{N}\right)
\]

Expandindo os logaritmos:

\[
\ln L_0 = N_1 \ln N_1 + N_0 \ln N_0 - N \ln N
\]

Essa expressão aparecia no denominador da caixa original e está correta
--- o erro estava apenas na fórmula geral do \(R^2\).

\begin{center}\rule{0.5\linewidth}{0.5pt}\end{center}

\subsubsection{4.1.3 Fórmula completa para uso
prático}\label{fuxf3rmula-completa-para-uso-pruxe1tico}

\[
\boxed{
R^2_{\text{McFadden}} = 1 - \frac{\ln L^*}{\,N_1 \ln N_1 + N_0 \ln N_0 - N \ln N\,}
}
\]

Ou, com a notação de logaritmo natural (\(\log_e\)):

\[
\boxed{
R^2_{\text{McFadden}} = 1 - \frac{L^*}{\,N_1\log_e N_1 + N_0\log_e N_0 - (N_1+N_0)\log_e(N_1+N_0)\,}
}
\]

\begin{center}\rule{0.5\linewidth}{0.5pt}\end{center}

\subsubsection{4.1.4 Exemplo numérico}\label{exemplo-numuxe9rico}

Dados: \(N_1 = 60\), \(N_0 = 40\), \(N = 100\).

Modelo nulo:\\
\(\ln L_0 = 60\ln(0{,}6) + 40\ln(0{,}4) \approx -67{,}3\)

Suponha que o modelo completo tenha \(\ln L^* = -50{,}0\):

\[
R^2_{\text{McFadden}} = 1 - \frac{-50{,}0}{-67{,}3} \approx 1 - 0{,}743 = 0{,}257
\]

Interpretação: o modelo completo reduz cerca de \textbf{25,7\%} da
incerteza (deviance) em relação ao modelo nulo --- análogo ao \(R^2\)
tradicional, mas na escala da log-verossimilhança.

\subsection{4.2 Acurácia e taxa de erro
aparente}\label{acuruxe1cia-e-taxa-de-erro-aparente}

\[\text{Acurácia} = \frac{\text{classificações corretas}}{n}, \qquad \text{Taxa de Erro} = 1 - \text{Acurácia}\]

A \textbf{taxa de erro aparente} é calculada no conjunto de treino e
tende a subestimar o erro real. O correto é calculá-la no
\textbf{conjunto de teste} (dados não vistos durante o ajuste).

\subsection{4.3 Matriz de confusão}\label{matriz-de-confusuxe3o}

Para um limiar de decisão \(\tau\) (geralmente \(0{,}5\)),
classifica-se:

\[\hat{y}_i = \begin{cases} 1 & \text{se } \hat{p}_i \geq \tau \\ 0 & \text{se } \hat{p}_i < \tau \end{cases}\]

A \textbf{matriz de confusão} cruza os valores reais com as previsões:

{\def\LTcaptype{none} % do not increment counter
\begin{longtable}[]{@{}
  >{\raggedright\arraybackslash}p{(\linewidth - 4\tabcolsep) * \real{0.2727}}
  >{\centering\arraybackslash}p{(\linewidth - 4\tabcolsep) * \real{0.3636}}
  >{\centering\arraybackslash}p{(\linewidth - 4\tabcolsep) * \real{0.3636}}@{}}
\toprule\noalign{}
\begin{minipage}[b]{\linewidth}\raggedright
\end{minipage} & \begin{minipage}[b]{\linewidth}\centering
\textbf{Predito 0}
\end{minipage} & \begin{minipage}[b]{\linewidth}\centering
\textbf{Predito 1}
\end{minipage} \\
\midrule\noalign{}
\endhead
\bottomrule\noalign{}
\endlastfoot
\textbf{Real 0} & Verdadeiro Negativo (VN) & Falso Positivo (FP) ---
Erro Tipo I \\
\textbf{Real 1} & Falso Negativo (FN) --- Erro Tipo II & Verdadeiro
Positivo (VP) \\
\end{longtable}
}

\textbf{Métricas derivadas:}

\[\text{Sensibilidade} = \frac{VP}{VP+FN} \qquad \text{Especificidade} = \frac{VN}{VN+FP}\]

\[\text{Precisão} = \frac{VP}{VP+FP} \qquad \text{F1-Score} = 2\cdot\frac{\text{Precisão} \times \text{Sensibilidade}}{\text{Precisão} + \text{Sensibilidade}}\]

\begin{quote}
\textbf{Escolha do limiar \(\tau\):} em contextos onde falsos negativos
são muito custosos (ex.: diagnóstico de câncer), adota-se
\(\tau < 0{,}5\) para aumentar a sensibilidade, aceitando mais falsos
positivos. A curva ROC auxilia nessa escolha.
\end{quote}

\subsection{4.4 Curva ROC e AUC}\label{curva-roc-e-auc}

A \textbf{Curva ROC} (\emph{Receiver Operating Characteristic}) plota a
\textbf{Sensibilidade (TPR)} versus \textbf{(1 − Especificidade) = FPR}
para \textbf{todos os possíveis limiares} \(\tau \in [0,1]\).

A \textbf{AUC} (\emph{Area Under the ROC Curve}) resume o desempenho:

{\def\LTcaptype{none} % do not increment counter
\begin{longtable}[]{@{}cl@{}}
\toprule\noalign{}
AUC & Interpretação \\
\midrule\noalign{}
\endhead
\bottomrule\noalign{}
\endlastfoot
\(0{,}5\) & Classificador aleatório (sem poder preditivo) \\
\(0{,}70 - 0{,}80\) & Desempenho aceitável \\
\(0{,}80 - 0{,}90\) & Bom desempenho \\
\(> 0{,}90\) & Desempenho excelente \\
\end{longtable}
}

A AUC também pode ser interpretada como a probabilidade de que, dado um
par aleatório (um positivo e um negativo), o modelo atribua maior
probabilidade ao positivo: \(\text{AUC} = P(\hat{p}_1 > \hat{p}_0)\).

\begin{center}\rule{0.5\linewidth}{0.5pt}\end{center}

\section{5. Testes de Hipóteses}\label{testes-de-hipuxf3teses}

\subsection{5.1 Teste da razão de verossimilhanças (TRV) --- modelo
global}\label{teste-da-razuxe3o-de-verossimilhanuxe7as-trv-modelo-global}

Avalia se o modelo ajustado é \textbf{globalmente} significativo em
relação ao modelo nulo (sem preditoras).

\textbf{Hipóteses:}
\[H_0: a_1 = a_2 = \cdots = a_p = 0 \qquad \text{vs.} \qquad H_1: \exists\, i: a_i \neq 0\]

\textbf{Estatística de teste:}
\[G = 2(L^* - L_0) \;\overset{H_0}{\longrightarrow}\; \chi^2_{(p)}\]

onde \(p\) é o número de preditoras (graus de liberdade). Rejeita-se
\(H_0\) se o valor-p \(< \alpha\).

\textbf{Protocolo de decisão:}

\begin{enumerate}
\def\labelenumi{\arabic{enumi}.}
\tightlist
\item
  Definir a população de interesse
\item
  Enunciar \(H_0\) e \(H_1\)
\item
  Escolher o TRV como teste
\item
  Fixar \(\alpha = 0{,}05\)
\item
  Calcular \(G = 2(L^* - L_0)\)
\item
  Obter valor-p: \(P(\chi^2_{(p)} \geq G)\)
\item
  Decisão: rejeitar \(H_0\) se valor-p \(< 0{,}05\); caso contrário, não
  rejeitar
\end{enumerate}

\subsection{5.2 Teste de Wald --- coeficientes
individuais}\label{teste-de-wald-coeficientes-individuais}

Avalia a significância de \textbf{cada coeficiente} separadamente.

\textbf{Hipóteses para o \(i\)-ésimo coeficiente:}
\[H_0: a_i = 0 \qquad \text{vs.} \qquad H_1: a_i \neq 0\]

\textbf{Estatística de Wald:}
\[W_i = \left(\frac{\hat{a}_i}{SE(\hat{a}_i)}\right)^2 \;\overset{H_0}{\longrightarrow}\; \chi^2_{(1)}\]

O erro padrão \(SE(\hat{a}_i)\) é a raiz quadrada do \(i\)-ésimo
elemento diagonal da matriz de covariância dos estimadores:

\[\widehat{\text{Var}}(\hat{\mathbf{a}}) = \left(\mathbf{X}^\top \widehat{W}\, \mathbf{X}\right)^{-1}, \quad \widehat{W} = \text{diag}\{\hat{y}_i(1-\hat{y}_i)\}\]

A estatística equivalente em escala normal é
\(z_i = \hat{a}_i / SE(\hat{a}_i) \sim \mathcal{N}(0,1)\), reportada
como ``z value'' na saída do R.

\textbf{Protocolo de decisão:}

\begin{enumerate}
\def\labelenumi{\arabic{enumi}.}
\tightlist
\item
  Definir a população
\item
  \(H_0: a_i = 0\) e \(H_1: a_i \neq 0\)
\item
  Escolher o Teste de Wald
\item
  Fixar \(\alpha = 0{,}05\)
\item
  Calcular \(W_i = \hat{a}_i^2 / SE(\hat{a}_i)^2\)
\item
  Obter valor-p: \(P(\chi^2_{(1)} \geq W_i)\)
\item
  Rejeitar \(H_0\) se valor-p \(< 0{,}05\)
\end{enumerate}

\begin{center}\rule{0.5\linewidth}{0.5pt}\end{center}

\section{6. Previsão com o Modelo}\label{previsuxe3o-com-o-modelo}

Com o modelo validado, a probabilidade predita para uma nova observação
\(\mathbf{x}^* = (x_1^*, \ldots, x_p^*)\) é:

\[\hat{p} = \frac{1}{1 + e^{-(\hat{a}_1 x_1^* + \cdots + \hat{a}_p x_p^* + \hat{b})}}\]

A classificação usa o limiar \(\tau\):

\[\hat{y} = \begin{cases} 1 & \text{se } \hat{p} \geq \tau \\ 0 & \text{se } \hat{p} < \tau \end{cases}\]

O limiar padrão \(\tau = 0{,}5\) pode ser ajustado com base na curva ROC
para equilibrar sensibilidade e especificidade conforme o contexto.

\begin{center}\rule{0.5\linewidth}{0.5pt}\end{center}

\section{7. Implementação em R}\label{implementauxe7uxe3o-em-r}

\begin{quote}
\textbf{Como usar:} ajuste as \textbf{6 variáveis} do bloco
\texttt{CONFIGURAÇÃO} abaixo para o seu arquivo \texttt{.csv}. Todo o
restante do código funciona automaticamente. Se o arquivo não existir,
dados simulados são gerados para demonstração.
\end{quote}

\subsection{7.1 Pacotes}\label{pacotes}

\begin{Shaded}
\begin{Highlighting}[]
\CommentTok{\# Instalação automática dos pacotes ausentes}
\NormalTok{pkgs\_necessarios }\OtherTok{\textless{}{-}} \FunctionTok{c}\NormalTok{(}\StringTok{"ggplot2"}\NormalTok{, }\StringTok{"dplyr"}\NormalTok{, }\StringTok{"pROC"}\NormalTok{, }\StringTok{"caret"}\NormalTok{,}
                      \StringTok{"DescTools"}\NormalTok{, }\StringTok{"gridExtra"}\NormalTok{, }\StringTok{"knitr"}\NormalTok{)}
\NormalTok{pkgs\_ausentes    }\OtherTok{\textless{}{-}}\NormalTok{ pkgs\_necessarios[}
  \SpecialCharTok{!}\NormalTok{pkgs\_necessarios }\SpecialCharTok{\%in\%} \FunctionTok{installed.packages}\NormalTok{()[, }\StringTok{"Package"}\NormalTok{]}
\NormalTok{]}
\ControlFlowTok{if}\NormalTok{ (}\FunctionTok{length}\NormalTok{(pkgs\_ausentes) }\SpecialCharTok{\textgreater{}} \DecValTok{0}\NormalTok{)}
  \FunctionTok{install.packages}\NormalTok{(pkgs\_ausentes, }\AttributeTok{repos =} \StringTok{"https://cran.r{-}project.org"}\NormalTok{,}
                   \AttributeTok{quiet =} \ConstantTok{TRUE}\NormalTok{)}

\FunctionTok{suppressPackageStartupMessages}\NormalTok{(\{}
  \FunctionTok{library}\NormalTok{(ggplot2)}
  \FunctionTok{library}\NormalTok{(dplyr)}
  \FunctionTok{library}\NormalTok{(pROC)}
  \FunctionTok{library}\NormalTok{(caret)}
  \FunctionTok{library}\NormalTok{(DescTools)}
  \FunctionTok{library}\NormalTok{(gridExtra)}
  \FunctionTok{library}\NormalTok{(knitr)}
\NormalTok{\})}
\end{Highlighting}
\end{Shaded}

\subsection{7.2 Configuração e importação dos
dados}\label{configurauxe7uxe3o-e-importauxe7uxe3o-dos-dados}

\begin{Shaded}
\begin{Highlighting}[]
\CommentTok{\# ════════════════════════════════════════════════════════════}
\CommentTok{\#  CONFIGURAÇÃO — ajuste apenas estas variáveis}
\CommentTok{\# ════════════════════════════════════════════════════════════}
\NormalTok{CAMINHO\_CSV       }\OtherTok{\textless{}{-}} \StringTok{"dados.csv"}      \CommentTok{\# caminho para o arquivo .csv}
\NormalTok{VARIAVEL\_RESPOSTA }\OtherTok{\textless{}{-}} \StringTok{"y"}              \CommentTok{\# coluna resposta (valores 0 e 1)}
\NormalTok{VARIAVEIS\_PRED    }\OtherTok{\textless{}{-}} \FunctionTok{c}\NormalTok{(}\StringTok{"x1"}\NormalTok{, }\StringTok{"x2"}\NormalTok{)   }\CommentTok{\# vetor com nomes das preditoras}
\NormalTok{LIMIAR\_DECISAO    }\OtherTok{\textless{}{-}} \FloatTok{0.5}             \CommentTok{\# limiar de classificação τ}
\NormalTok{PROPORCAO\_TREINO  }\OtherTok{\textless{}{-}} \FloatTok{0.70}            \CommentTok{\# fração para treino}
\NormalTok{SEMENTE           }\OtherTok{\textless{}{-}} \DecValTok{42}              \CommentTok{\# semente para reprodutibilidade}
\CommentTok{\# ════════════════════════════════════════════════════════════}

\CommentTok{\# ── Geração de dados simulados (apenas se o .csv não existir) ──────────────}
\ControlFlowTok{if}\NormalTok{ (}\SpecialCharTok{!}\FunctionTok{file.exists}\NormalTok{(CAMINHO\_CSV)) \{}
  \FunctionTok{message}\NormalTok{(}\StringTok{"Arquivo não encontrado — gerando dados simulados."}\NormalTok{)}
  \FunctionTok{set.seed}\NormalTok{(SEMENTE)}
\NormalTok{  n\_sim  }\OtherTok{\textless{}{-}} \DecValTok{350}
\NormalTok{  x1\_sim }\OtherTok{\textless{}{-}} \FunctionTok{round}\NormalTok{(}\FunctionTok{runif}\NormalTok{(n\_sim, }\DecValTok{18}\NormalTok{, }\DecValTok{65}\NormalTok{), }\DecValTok{1}\NormalTok{)   }\CommentTok{\# ex.: idade}
\NormalTok{  x2\_sim }\OtherTok{\textless{}{-}} \FunctionTok{round}\NormalTok{(}\FunctionTok{rnorm}\NormalTok{(n\_sim, }\DecValTok{50}\NormalTok{, }\DecValTok{10}\NormalTok{),  }\DecValTok{1}\NormalTok{)   }\CommentTok{\# ex.: pontuação}
\NormalTok{  z\_sim  }\OtherTok{\textless{}{-}} \FloatTok{0.07} \SpecialCharTok{*}\NormalTok{ x1\_sim }\SpecialCharTok{+} \FloatTok{0.06} \SpecialCharTok{*}\NormalTok{ x2\_sim }\SpecialCharTok{{-}} \FloatTok{7.0}
\NormalTok{  p\_sim  }\OtherTok{\textless{}{-}} \DecValTok{1} \SpecialCharTok{/}\NormalTok{ (}\DecValTok{1} \SpecialCharTok{+} \FunctionTok{exp}\NormalTok{(}\SpecialCharTok{{-}}\NormalTok{z\_sim))}
\NormalTok{  y\_sim  }\OtherTok{\textless{}{-}} \FunctionTok{rbinom}\NormalTok{(n\_sim, }\DecValTok{1}\NormalTok{, p\_sim)}
\NormalTok{  df\_sim }\OtherTok{\textless{}{-}} \FunctionTok{data.frame}\NormalTok{(}\AttributeTok{y =}\NormalTok{ y\_sim, }\AttributeTok{x1 =}\NormalTok{ x1\_sim, }\AttributeTok{x2 =}\NormalTok{ x2\_sim)}
  \FunctionTok{write.csv}\NormalTok{(df\_sim, CAMINHO\_CSV, }\AttributeTok{row.names =} \ConstantTok{FALSE}\NormalTok{)}
  \FunctionTok{message}\NormalTok{(}\FunctionTok{sprintf}\NormalTok{(}\StringTok{"Arquivo \textquotesingle{}\%s\textquotesingle{} criado com \%d observações."}\NormalTok{, CAMINHO\_CSV, n\_sim))}
\NormalTok{\}}

\CommentTok{\# ── Leitura ─────────────────────────────────────────────────────────────────}
\NormalTok{raw }\OtherTok{\textless{}{-}} \FunctionTok{read.csv}\NormalTok{(CAMINHO\_CSV, }\AttributeTok{stringsAsFactors =} \ConstantTok{FALSE}\NormalTok{)}
\FunctionTok{cat}\NormalTok{(}\FunctionTok{sprintf}\NormalTok{(}\StringTok{"Dataset lido: \%d linhas × \%d colunas}\SpecialCharTok{\textbackslash{}n}\StringTok{"}\NormalTok{, }\FunctionTok{nrow}\NormalTok{(raw), }\FunctionTok{ncol}\NormalTok{(raw)))}
\FunctionTok{cat}\NormalTok{(}\StringTok{"Colunas:"}\NormalTok{, }\FunctionTok{paste}\NormalTok{(}\FunctionTok{names}\NormalTok{(raw), }\AttributeTok{collapse =} \StringTok{", "}\NormalTok{), }\StringTok{"}\SpecialCharTok{\textbackslash{}n}\StringTok{"}\NormalTok{)}

\CommentTok{\# ── Validação das colunas configuradas ──────────────────────────────────────}
\NormalTok{cols\_req     }\OtherTok{\textless{}{-}} \FunctionTok{c}\NormalTok{(VARIAVEL\_RESPOSTA, VARIAVEIS\_PRED)}
\NormalTok{cols\_missing }\OtherTok{\textless{}{-}} \FunctionTok{setdiff}\NormalTok{(cols\_req, }\FunctionTok{names}\NormalTok{(raw))}
\ControlFlowTok{if}\NormalTok{ (}\FunctionTok{length}\NormalTok{(cols\_missing) }\SpecialCharTok{\textgreater{}} \DecValTok{0}\NormalTok{)}
  \FunctionTok{stop}\NormalTok{(}\StringTok{"Colunas ausentes no CSV: "}\NormalTok{, }\FunctionTok{paste}\NormalTok{(cols\_missing, }\AttributeTok{collapse =} \StringTok{", "}\NormalTok{),}
       \StringTok{"}\SpecialCharTok{\textbackslash{}n}\StringTok{Ajuste VARIAVEL\_RESPOSTA e VARIAVEIS\_PRED."}\NormalTok{)}

\CommentTok{\# ── Data frame de trabalho ───────────────────────────────────────────────────}
\NormalTok{dados           }\OtherTok{\textless{}{-}}\NormalTok{ raw[, cols\_req]}
\FunctionTok{names}\NormalTok{(dados)[}\DecValTok{1}\NormalTok{] }\OtherTok{\textless{}{-}} \StringTok{"y"}
\NormalTok{dados}\SpecialCharTok{$}\NormalTok{y         }\OtherTok{\textless{}{-}} \FunctionTok{as.integer}\NormalTok{(}\FunctionTok{as.character}\NormalTok{(dados}\SpecialCharTok{$}\NormalTok{y))}

\CommentTok{\# Remoção de NAs}
\NormalTok{n\_antes  }\OtherTok{\textless{}{-}} \FunctionTok{nrow}\NormalTok{(dados)}
\NormalTok{dados    }\OtherTok{\textless{}{-}} \FunctionTok{na.omit}\NormalTok{(dados)}
\NormalTok{n\_depois }\OtherTok{\textless{}{-}} \FunctionTok{nrow}\NormalTok{(dados)}
\ControlFlowTok{if}\NormalTok{ (n\_antes }\SpecialCharTok{\textgreater{}}\NormalTok{ n\_depois)}
  \FunctionTok{cat}\NormalTok{(}\FunctionTok{sprintf}\NormalTok{(}\StringTok{"Removidas \%d linha(s) com NA.}\SpecialCharTok{\textbackslash{}n}\StringTok{"}\NormalTok{, n\_antes }\SpecialCharTok{{-}}\NormalTok{ n\_depois))}

\FunctionTok{cat}\NormalTok{(}\FunctionTok{sprintf}\NormalTok{(}\StringTok{"Observações válidas: \%d}\SpecialCharTok{\textbackslash{}n}\StringTok{"}\NormalTok{, n\_depois))}
\end{Highlighting}
\end{Shaded}

\subsection{7.3 Análise exploratória}\label{anuxe1lise-exploratuxf3ria}

\begin{Shaded}
\begin{Highlighting}[]
\NormalTok{N1 }\OtherTok{\textless{}{-}} \FunctionTok{sum}\NormalTok{(dados}\SpecialCharTok{$}\NormalTok{y }\SpecialCharTok{==} \DecValTok{1}\NormalTok{)}
\NormalTok{N0 }\OtherTok{\textless{}{-}} \FunctionTok{sum}\NormalTok{(dados}\SpecialCharTok{$}\NormalTok{y }\SpecialCharTok{==} \DecValTok{0}\NormalTok{)}
\NormalTok{N  }\OtherTok{\textless{}{-}} \FunctionTok{nrow}\NormalTok{(dados)}

\FunctionTok{cat}\NormalTok{(}\FunctionTok{sprintf}\NormalTok{(}
  \StringTok{"Classe 1 (sucesso):  \%d  (\%.1f\%\%)}\SpecialCharTok{\textbackslash{}n}\StringTok{Classe 0 (fracasso): \%d  (\%.1f\%\%)}\SpecialCharTok{\textbackslash{}n}\StringTok{"}\NormalTok{,}
\NormalTok{  N1, }\DecValTok{100} \SpecialCharTok{*}\NormalTok{ N1 }\SpecialCharTok{/}\NormalTok{ N, N0, }\DecValTok{100} \SpecialCharTok{*}\NormalTok{ N0 }\SpecialCharTok{/}\NormalTok{ N}
\NormalTok{))}

\NormalTok{dados}\SpecialCharTok{$}\NormalTok{y\_fator }\OtherTok{\textless{}{-}} \FunctionTok{factor}\NormalTok{(dados}\SpecialCharTok{$}\NormalTok{y, }\AttributeTok{levels =} \FunctionTok{c}\NormalTok{(}\DecValTok{0}\NormalTok{, }\DecValTok{1}\NormalTok{),}
                         \AttributeTok{labels =} \FunctionTok{c}\NormalTok{(}\StringTok{"Fracasso (0)"}\NormalTok{, }\StringTok{"Sucesso (1)"}\NormalTok{))}
\NormalTok{pal }\OtherTok{\textless{}{-}} \FunctionTok{c}\NormalTok{(}\StringTok{"Fracasso (0)"} \OtherTok{=} \StringTok{"\#C0392B"}\NormalTok{, }\StringTok{"Sucesso (1)"} \OtherTok{=} \StringTok{"\#27AE60"}\NormalTok{)}

\CommentTok{\# Gráfico de dispersão}
\NormalTok{p\_disp }\OtherTok{\textless{}{-}} \FunctionTok{ggplot}\NormalTok{(dados,}
                 \FunctionTok{aes\_string}\NormalTok{(}\AttributeTok{x =}\NormalTok{ VARIAVEIS\_PRED[}\DecValTok{1}\NormalTok{],}
                             \AttributeTok{y =} \ControlFlowTok{if}\NormalTok{ (}\FunctionTok{length}\NormalTok{(VARIAVEIS\_PRED) }\SpecialCharTok{\textgreater{}} \DecValTok{1}\NormalTok{) VARIAVEIS\_PRED[}\DecValTok{2}\NormalTok{]}
                                 \ControlFlowTok{else}\NormalTok{ VARIAVEIS\_PRED[}\DecValTok{1}\NormalTok{],}
                             \AttributeTok{color =} \StringTok{"y\_fator"}\NormalTok{, }\AttributeTok{shape =} \StringTok{"y\_fator"}\NormalTok{)) }\SpecialCharTok{+}
  \FunctionTok{geom\_jitter}\NormalTok{(}\AttributeTok{alpha =} \FloatTok{0.65}\NormalTok{, }\AttributeTok{size =} \FloatTok{1.8}\NormalTok{, }\AttributeTok{width =} \FloatTok{0.25}\NormalTok{, }\AttributeTok{height =} \FloatTok{0.25}\NormalTok{) }\SpecialCharTok{+}
  \FunctionTok{scale\_color\_manual}\NormalTok{(}\AttributeTok{values =}\NormalTok{ pal) }\SpecialCharTok{+}
  \FunctionTok{scale\_shape\_manual}\NormalTok{(}\AttributeTok{values =} \FunctionTok{c}\NormalTok{(}\StringTok{"Fracasso (0)"} \OtherTok{=} \DecValTok{16}\NormalTok{, }\StringTok{"Sucesso (1)"} \OtherTok{=} \DecValTok{17}\NormalTok{)) }\SpecialCharTok{+}
  \FunctionTok{labs}\NormalTok{(}\AttributeTok{title =} \StringTok{"Dispersão por Classe"}\NormalTok{,}
       \AttributeTok{color =} \ConstantTok{NULL}\NormalTok{, }\AttributeTok{shape =} \ConstantTok{NULL}\NormalTok{) }\SpecialCharTok{+}
  \FunctionTok{theme\_minimal}\NormalTok{(}\AttributeTok{base\_size =} \DecValTok{11}\NormalTok{) }\SpecialCharTok{+}
  \FunctionTok{theme}\NormalTok{(}\AttributeTok{legend.position =} \StringTok{"bottom"}\NormalTok{)}

\CommentTok{\# Boxplots por preditora}
\NormalTok{plots\_bp }\OtherTok{\textless{}{-}} \FunctionTok{lapply}\NormalTok{(VARIAVEIS\_PRED, }\ControlFlowTok{function}\NormalTok{(v) \{}
  \FunctionTok{ggplot}\NormalTok{(dados, }\FunctionTok{aes\_string}\NormalTok{(}\AttributeTok{x =} \StringTok{"y\_fator"}\NormalTok{, }\AttributeTok{y =}\NormalTok{ v, }\AttributeTok{fill =} \StringTok{"y\_fator"}\NormalTok{)) }\SpecialCharTok{+}
    \FunctionTok{geom\_boxplot}\NormalTok{(}\AttributeTok{alpha =} \FloatTok{0.75}\NormalTok{, }\AttributeTok{outlier.shape =} \DecValTok{21}\NormalTok{,}
                 \AttributeTok{outlier.fill =} \StringTok{"white"}\NormalTok{, }\AttributeTok{outlier.stroke =} \FloatTok{0.5}\NormalTok{,}
                 \AttributeTok{outlier.size =} \FloatTok{1.5}\NormalTok{) }\SpecialCharTok{+}
    \FunctionTok{scale\_fill\_manual}\NormalTok{(}\AttributeTok{values =}\NormalTok{ pal) }\SpecialCharTok{+}
    \FunctionTok{labs}\NormalTok{(}\AttributeTok{x =} \ConstantTok{NULL}\NormalTok{, }\AttributeTok{y =}\NormalTok{ v,}
         \AttributeTok{title =} \FunctionTok{paste}\NormalTok{(}\StringTok{"Distribuição de"}\NormalTok{, v)) }\SpecialCharTok{+}
    \FunctionTok{theme\_minimal}\NormalTok{(}\AttributeTok{base\_size =} \DecValTok{11}\NormalTok{) }\SpecialCharTok{+}
    \FunctionTok{theme}\NormalTok{(}\AttributeTok{legend.position =} \StringTok{"none"}\NormalTok{,}
          \AttributeTok{axis.text.x =} \FunctionTok{element\_text}\NormalTok{(}\AttributeTok{size =} \DecValTok{9}\NormalTok{))}
\NormalTok{\})}

\FunctionTok{do.call}\NormalTok{(grid.arrange, }\FunctionTok{c}\NormalTok{(}\FunctionTok{list}\NormalTok{(p\_disp), plots\_bp, }\FunctionTok{list}\NormalTok{(}\AttributeTok{ncol =} \DecValTok{3}\NormalTok{)))}
\end{Highlighting}
\end{Shaded}

\subsection{7.4 Divisão treino/teste e ajuste do
modelo}\label{divisuxe3o-treinoteste-e-ajuste-do-modelo}

\begin{Shaded}
\begin{Highlighting}[]
\FunctionTok{set.seed}\NormalTok{(SEMENTE)}

\CommentTok{\# Partição estratificada (mantém proporção de classes)}
\NormalTok{idx\_treino }\OtherTok{\textless{}{-}} \FunctionTok{createDataPartition}\NormalTok{(dados}\SpecialCharTok{$}\NormalTok{y, }\AttributeTok{p =}\NormalTok{ PROPORCAO\_TREINO, }\AttributeTok{list =} \ConstantTok{FALSE}\NormalTok{)}
\NormalTok{treino     }\OtherTok{\textless{}{-}}\NormalTok{ dados[ idx\_treino, ]}
\NormalTok{teste      }\OtherTok{\textless{}{-}}\NormalTok{ dados[}\SpecialCharTok{{-}}\NormalTok{idx\_treino, ]}

\FunctionTok{cat}\NormalTok{(}\FunctionTok{sprintf}\NormalTok{(}\StringTok{"Treino: \%d obs.  |  Teste: \%d obs.}\SpecialCharTok{\textbackslash{}n}\StringTok{"}\NormalTok{, }\FunctionTok{nrow}\NormalTok{(treino), }\FunctionTok{nrow}\NormalTok{(teste)))}

\CommentTok{\# Fórmula do modelo}
\NormalTok{formula\_modelo }\OtherTok{\textless{}{-}} \FunctionTok{as.formula}\NormalTok{(}
  \FunctionTok{paste}\NormalTok{(}\StringTok{"y \textasciitilde{}"}\NormalTok{, }\FunctionTok{paste}\NormalTok{(VARIAVEIS\_PRED, }\AttributeTok{collapse =} \StringTok{" + "}\NormalTok{))}
\NormalTok{)}

\CommentTok{\# Ajuste via GLM — família binomial, ligação logit (MV via IRLS)}
\NormalTok{modelo }\OtherTok{\textless{}{-}} \FunctionTok{glm}\NormalTok{(formula\_modelo, }\AttributeTok{data =}\NormalTok{ treino,}
              \AttributeTok{family =} \FunctionTok{binomial}\NormalTok{(}\AttributeTok{link =} \StringTok{"logit"}\NormalTok{))}

\FunctionTok{print}\NormalTok{(}\FunctionTok{summary}\NormalTok{(modelo))}
\end{Highlighting}
\end{Shaded}

\subsection{7.5 Coeficientes e odds
ratios}\label{coeficientes-e-odds-ratios}

\begin{Shaded}
\begin{Highlighting}[]
\NormalTok{ci\_logit  }\OtherTok{\textless{}{-}} \FunctionTok{suppressMessages}\NormalTok{(}\FunctionTok{confint}\NormalTok{(modelo))}
\NormalTok{or\_table  }\OtherTok{\textless{}{-}} \FunctionTok{exp}\NormalTok{(}\FunctionTok{cbind}\NormalTok{(}\AttributeTok{OR =} \FunctionTok{coef}\NormalTok{(modelo), ci\_logit))}
\FunctionTok{colnames}\NormalTok{(or\_table) }\OtherTok{\textless{}{-}} \FunctionTok{c}\NormalTok{(}\StringTok{"OR"}\NormalTok{, }\StringTok{"IC 2,5\%"}\NormalTok{, }\StringTok{"IC 97,5\%"}\NormalTok{)}

\FunctionTok{cat}\NormalTok{(}\StringTok{"=== Odds Ratios com Intervalo de Confiança 95\% ===}\SpecialCharTok{\textbackslash{}n}\StringTok{"}\NormalTok{)}
\FunctionTok{kable}\NormalTok{(}\FunctionTok{round}\NormalTok{(or\_table, }\DecValTok{4}\NormalTok{),}
      \AttributeTok{caption =} \StringTok{"Tabela 1. Coeficientes, Odds Ratios e IC 95\%"}\NormalTok{)}
\end{Highlighting}
\end{Shaded}

\subsection{7.6 Pseudo-R² de McFadden}\label{pseudo-ruxb2-de-mcfadden-1}

\begin{Shaded}
\begin{Highlighting}[]
\NormalTok{L\_star    }\OtherTok{\textless{}{-}} \FunctionTok{as.numeric}\NormalTok{(}\FunctionTok{logLik}\NormalTok{(modelo))}
\NormalTok{L\_nulo\_val }\OtherTok{\textless{}{-}}\NormalTok{ N1 }\SpecialCharTok{*} \FunctionTok{log}\NormalTok{(N1) }\SpecialCharTok{+}\NormalTok{ N0 }\SpecialCharTok{*} \FunctionTok{log}\NormalTok{(N0) }\SpecialCharTok{{-}}\NormalTok{ N }\SpecialCharTok{*} \FunctionTok{log}\NormalTok{(N)}
\NormalTok{r2\_mcf    }\OtherTok{\textless{}{-}} \SpecialCharTok{{-}}\NormalTok{L\_star }\SpecialCharTok{/}\NormalTok{ L\_nulo\_val}

\FunctionTok{cat}\NormalTok{(}\FunctionTok{sprintf}\NormalTok{(}\StringTok{"Log{-}verossimilhança do modelo (L*): \%.4f}\SpecialCharTok{\textbackslash{}n}\StringTok{"}\NormalTok{, L\_star))}
\FunctionTok{cat}\NormalTok{(}\FunctionTok{sprintf}\NormalTok{(}\StringTok{"Log{-}verossimilhança do nulo  (L0): \%.4f}\SpecialCharTok{\textbackslash{}n}\StringTok{"}\NormalTok{, L\_nulo\_val))}
\FunctionTok{cat}\NormalTok{(}\FunctionTok{sprintf}\NormalTok{(}\StringTok{"Pseudo{-}R² de McFadden:              \%.4f}\SpecialCharTok{\textbackslash{}n}\StringTok{"}\NormalTok{, r2\_mcf))}
\FunctionTok{cat}\NormalTok{(}\StringTok{"Avaliação:"}\NormalTok{,}
\NormalTok{    dplyr}\SpecialCharTok{::}\FunctionTok{case\_when}\NormalTok{(}
\NormalTok{      r2\_mcf }\SpecialCharTok{\textless{}} \FloatTok{0.20} \SpecialCharTok{\textasciitilde{}} \StringTok{"Ajuste fraco (\textless{} 0,20)"}\NormalTok{,}
\NormalTok{      r2\_mcf }\SpecialCharTok{\textless{}} \FloatTok{0.40} \SpecialCharTok{\textasciitilde{}} \StringTok{"Bom ajuste (0,20 – 0,40)"}\NormalTok{,}
      \ConstantTok{TRUE}          \SpecialCharTok{\textasciitilde{}} \StringTok{"Ajuste muito bom (\textgreater{} 0,40)"}
\NormalTok{    ), }\StringTok{"}\SpecialCharTok{\textbackslash{}n}\StringTok{"}\NormalTok{)}
\end{Highlighting}
\end{Shaded}

\subsection{7.7 Teste da razão de
verossimilhanças}\label{teste-da-razuxe3o-de-verossimilhanuxe7as}

\begin{Shaded}
\begin{Highlighting}[]
\NormalTok{modelo\_nulo }\OtherTok{\textless{}{-}} \FunctionTok{glm}\NormalTok{(y }\SpecialCharTok{\textasciitilde{}} \DecValTok{1}\NormalTok{, }\AttributeTok{data =}\NormalTok{ treino, }\AttributeTok{family =} \FunctionTok{binomial}\NormalTok{(}\AttributeTok{link =} \StringTok{"logit"}\NormalTok{))}

\NormalTok{G        }\OtherTok{\textless{}{-}} \DecValTok{2} \SpecialCharTok{*}\NormalTok{ (}\FunctionTok{as.numeric}\NormalTok{(}\FunctionTok{logLik}\NormalTok{(modelo)) }\SpecialCharTok{{-}} \FunctionTok{as.numeric}\NormalTok{(}\FunctionTok{logLik}\NormalTok{(modelo\_nulo)))}
\NormalTok{gl\_trv   }\OtherTok{\textless{}{-}} \FunctionTok{length}\NormalTok{(VARIAVEIS\_PRED)}
\NormalTok{pval\_trv }\OtherTok{\textless{}{-}} \DecValTok{1} \SpecialCharTok{{-}} \FunctionTok{pchisq}\NormalTok{(G, }\AttributeTok{df =}\NormalTok{ gl\_trv)}

\FunctionTok{cat}\NormalTok{(}\StringTok{"H0: a1 = a2 = ... = ap = 0}\SpecialCharTok{\textbackslash{}n}\StringTok{"}\NormalTok{)}
\FunctionTok{cat}\NormalTok{(}\StringTok{"H1: pelo menos um ai ≠ 0}\SpecialCharTok{\textbackslash{}n\textbackslash{}n}\StringTok{"}\NormalTok{)}
\FunctionTok{cat}\NormalTok{(}\FunctionTok{sprintf}\NormalTok{(}\StringTok{"Estatística G:      \%.4f}\SpecialCharTok{\textbackslash{}n}\StringTok{"}\NormalTok{, G))}
\FunctionTok{cat}\NormalTok{(}\FunctionTok{sprintf}\NormalTok{(}\StringTok{"Graus de liberdade: \%d}\SpecialCharTok{\textbackslash{}n}\StringTok{"}\NormalTok{,   gl\_trv))}
\FunctionTok{cat}\NormalTok{(}\FunctionTok{sprintf}\NormalTok{(}\StringTok{"Valor{-}p:            \%.2e}\SpecialCharTok{\textbackslash{}n}\StringTok{"}\NormalTok{, pval\_trv))}
\FunctionTok{cat}\NormalTok{(}\StringTok{"Decisão (α = 0,05):"}\NormalTok{,}
    \FunctionTok{ifelse}\NormalTok{(pval\_trv }\SpecialCharTok{\textless{}} \FloatTok{0.05}\NormalTok{,}
           \StringTok{"Rejeitar H0 — modelo globalmente significativo."}\NormalTok{,}
           \StringTok{"Não rejeitar H0."}\NormalTok{), }\StringTok{"}\SpecialCharTok{\textbackslash{}n}\StringTok{"}\NormalTok{)}
\end{Highlighting}
\end{Shaded}

\subsection{7.8 Teste de Wald}\label{teste-de-wald}

\begin{Shaded}
\begin{Highlighting}[]
\NormalTok{coefs\_v   }\OtherTok{\textless{}{-}} \FunctionTok{coef}\NormalTok{(modelo)}
\NormalTok{ep\_v      }\OtherTok{\textless{}{-}} \FunctionTok{sqrt}\NormalTok{(}\FunctionTok{diag}\NormalTok{(}\FunctionTok{vcov}\NormalTok{(modelo)))}
\NormalTok{w\_stat\_v  }\OtherTok{\textless{}{-}}\NormalTok{ (coefs\_v }\SpecialCharTok{/}\NormalTok{ ep\_v)}\SpecialCharTok{\^{}}\DecValTok{2}
\NormalTok{p\_wald\_v  }\OtherTok{\textless{}{-}} \DecValTok{1} \SpecialCharTok{{-}} \FunctionTok{pchisq}\NormalTok{(w\_stat\_v, }\AttributeTok{df =} \DecValTok{1}\NormalTok{)}

\NormalTok{wald\_df }\OtherTok{\textless{}{-}} \FunctionTok{data.frame}\NormalTok{(}
  \AttributeTok{Coeficiente   =} \FunctionTok{round}\NormalTok{(coefs\_v,   }\DecValTok{4}\NormalTok{),}
  \AttributeTok{Erro\_Padrao   =} \FunctionTok{round}\NormalTok{(ep\_v,      }\DecValTok{4}\NormalTok{),}
  \AttributeTok{Estatistica\_W =} \FunctionTok{round}\NormalTok{(w\_stat\_v,  }\DecValTok{4}\NormalTok{),}
  \AttributeTok{Valor\_p       =} \FunctionTok{round}\NormalTok{(p\_wald\_v,  }\DecValTok{6}\NormalTok{),}
  \AttributeTok{Decisao       =} \FunctionTok{ifelse}\NormalTok{(p\_wald\_v }\SpecialCharTok{\textless{}} \FloatTok{0.05}\NormalTok{, }\StringTok{"Rejeitar H0 (*)"}\NormalTok{, }\StringTok{"Não rejeitar H0"}\NormalTok{)}
\NormalTok{)}

\FunctionTok{cat}\NormalTok{(}\StringTok{"H0: a\_i = 0  |  H1: a\_i ≠ 0}\SpecialCharTok{\textbackslash{}n\textbackslash{}n}\StringTok{"}\NormalTok{)}
\FunctionTok{kable}\NormalTok{(wald\_df, }\AttributeTok{caption =} \StringTok{"Tabela 2. Teste de Wald para cada coeficiente"}\NormalTok{)}
\end{Highlighting}
\end{Shaded}

\subsection{7.9 Previsão e avaliação no conjunto de
teste}\label{previsuxe3o-e-avaliauxe7uxe3o-no-conjunto-de-teste}

\begin{Shaded}
\begin{Highlighting}[]
\NormalTok{prob\_teste }\OtherTok{\textless{}{-}} \FunctionTok{predict}\NormalTok{(modelo, }\AttributeTok{newdata =}\NormalTok{ teste, }\AttributeTok{type =} \StringTok{"response"}\NormalTok{)}
\NormalTok{y\_pred\_r   }\OtherTok{\textless{}{-}} \FunctionTok{ifelse}\NormalTok{(prob\_teste }\SpecialCharTok{\textgreater{}=}\NormalTok{ LIMIAR\_DECISAO, }\DecValTok{1}\NormalTok{, }\DecValTok{0}\NormalTok{)}

\NormalTok{acuracia\_r  }\OtherTok{\textless{}{-}} \FunctionTok{mean}\NormalTok{(y\_pred\_r }\SpecialCharTok{==}\NormalTok{ teste}\SpecialCharTok{$}\NormalTok{y)}
\NormalTok{taxa\_erro\_r }\OtherTok{\textless{}{-}} \DecValTok{1} \SpecialCharTok{{-}}\NormalTok{ acuracia\_r}

\FunctionTok{cat}\NormalTok{(}\FunctionTok{sprintf}\NormalTok{(}\StringTok{"Limiar de decisão (τ): \%.2f}\SpecialCharTok{\textbackslash{}n}\StringTok{"}\NormalTok{, LIMIAR\_DECISAO))}
\FunctionTok{cat}\NormalTok{(}\FunctionTok{sprintf}\NormalTok{(}\StringTok{"Acurácia:              \%.4f  (\%.1f\%\%)}\SpecialCharTok{\textbackslash{}n}\StringTok{"}\NormalTok{,}
\NormalTok{            acuracia\_r, }\DecValTok{100} \SpecialCharTok{*}\NormalTok{ acuracia\_r))}
\FunctionTok{cat}\NormalTok{(}\FunctionTok{sprintf}\NormalTok{(}\StringTok{"Taxa de erro:          \%.4f  (\%.1f\%\%)}\SpecialCharTok{\textbackslash{}n\textbackslash{}n}\StringTok{"}\NormalTok{,}
\NormalTok{            taxa\_erro\_r, }\DecValTok{100} \SpecialCharTok{*}\NormalTok{ taxa\_erro\_r))}

\NormalTok{cm\_r }\OtherTok{\textless{}{-}} \FunctionTok{confusionMatrix}\NormalTok{(}
  \FunctionTok{factor}\NormalTok{(y\_pred\_r, }\AttributeTok{levels =} \FunctionTok{c}\NormalTok{(}\DecValTok{0}\NormalTok{, }\DecValTok{1}\NormalTok{)),}
  \FunctionTok{factor}\NormalTok{(teste}\SpecialCharTok{$}\NormalTok{y,  }\AttributeTok{levels =} \FunctionTok{c}\NormalTok{(}\DecValTok{0}\NormalTok{, }\DecValTok{1}\NormalTok{)),}
  \AttributeTok{positive =} \StringTok{"1"}
\NormalTok{)}
\FunctionTok{print}\NormalTok{(cm\_r)}
\end{Highlighting}
\end{Shaded}

\subsection{7.10 Visualizações: matriz de confusão e curva
ROC}\label{visualizauxe7uxf5es-matriz-de-confusuxe3o-e-curva-roc}

\begin{Shaded}
\begin{Highlighting}[]
\CommentTok{\# ── Matriz de confusão ──────────────────────────────────────────────────────}
\NormalTok{cm\_df\_r }\OtherTok{\textless{}{-}} \FunctionTok{as.data.frame}\NormalTok{(cm\_r}\SpecialCharTok{$}\NormalTok{table)}
\FunctionTok{names}\NormalTok{(cm\_df\_r) }\OtherTok{\textless{}{-}} \FunctionTok{c}\NormalTok{(}\StringTok{"Predito"}\NormalTok{, }\StringTok{"Real"}\NormalTok{, }\StringTok{"N"}\NormalTok{)}

\NormalTok{p\_mc }\OtherTok{\textless{}{-}} \FunctionTok{ggplot}\NormalTok{(cm\_df\_r, }\FunctionTok{aes}\NormalTok{(}\AttributeTok{x =}\NormalTok{ Real, }\AttributeTok{y =}\NormalTok{ Predito, }\AttributeTok{fill =}\NormalTok{ N)) }\SpecialCharTok{+}
  \FunctionTok{geom\_tile}\NormalTok{(}\AttributeTok{color =} \StringTok{"white"}\NormalTok{, }\AttributeTok{size =} \FloatTok{1.2}\NormalTok{) }\SpecialCharTok{+}
  \FunctionTok{geom\_text}\NormalTok{(}\FunctionTok{aes}\NormalTok{(}\AttributeTok{label =}\NormalTok{ N), }\AttributeTok{size =} \DecValTok{9}\NormalTok{, }\AttributeTok{fontface =} \StringTok{"bold"}\NormalTok{, }\AttributeTok{color =} \StringTok{"white"}\NormalTok{) }\SpecialCharTok{+}
  \FunctionTok{scale\_fill\_gradient}\NormalTok{(}\AttributeTok{low =} \StringTok{"\#AED6F1"}\NormalTok{, }\AttributeTok{high =} \StringTok{"\#1A5276"}\NormalTok{) }\SpecialCharTok{+}
  \FunctionTok{scale\_x\_discrete}\NormalTok{(}\AttributeTok{labels =} \FunctionTok{c}\NormalTok{(}\StringTok{"0"} \OtherTok{=} \StringTok{"Real 0"}\NormalTok{, }\StringTok{"1"} \OtherTok{=} \StringTok{"Real 1"}\NormalTok{)) }\SpecialCharTok{+}
  \FunctionTok{scale\_y\_discrete}\NormalTok{(}\AttributeTok{labels =} \FunctionTok{c}\NormalTok{(}\StringTok{"0"} \OtherTok{=} \StringTok{"Pred 0"}\NormalTok{, }\StringTok{"1"} \OtherTok{=} \StringTok{"Pred 1"}\NormalTok{)) }\SpecialCharTok{+}
  \FunctionTok{labs}\NormalTok{(}\AttributeTok{title =} \StringTok{"Matriz de Confusão"}\NormalTok{,}
       \AttributeTok{x =} \StringTok{"Valor Real"}\NormalTok{, }\AttributeTok{y =} \StringTok{"Valor Predito"}\NormalTok{, }\AttributeTok{fill =} \StringTok{"n"}\NormalTok{) }\SpecialCharTok{+}
  \FunctionTok{theme\_minimal}\NormalTok{(}\AttributeTok{base\_size =} \DecValTok{12}\NormalTok{) }\SpecialCharTok{+}
  \FunctionTok{theme}\NormalTok{(}\AttributeTok{panel.grid =} \FunctionTok{element\_blank}\NormalTok{(),}
        \AttributeTok{axis.text  =} \FunctionTok{element\_text}\NormalTok{(}\AttributeTok{size =} \DecValTok{11}\NormalTok{))}

\CommentTok{\# ── Curva ROC ───────────────────────────────────────────────────────────────}
\NormalTok{roc\_r   }\OtherTok{\textless{}{-}} \FunctionTok{roc}\NormalTok{(teste}\SpecialCharTok{$}\NormalTok{y, prob\_teste, }\AttributeTok{quiet =} \ConstantTok{TRUE}\NormalTok{)}
\NormalTok{auc\_r   }\OtherTok{\textless{}{-}} \FunctionTok{auc}\NormalTok{(roc\_r)}
\NormalTok{roc\_df\_r }\OtherTok{\textless{}{-}} \FunctionTok{data.frame}\NormalTok{(}
  \AttributeTok{FPR =} \DecValTok{1} \SpecialCharTok{{-}}\NormalTok{ roc\_r}\SpecialCharTok{$}\NormalTok{specificities,}
  \AttributeTok{TPR =}\NormalTok{     roc\_r}\SpecialCharTok{$}\NormalTok{sensitivities}
\NormalTok{)}

\NormalTok{p\_roc }\OtherTok{\textless{}{-}} \FunctionTok{ggplot}\NormalTok{(roc\_df\_r, }\FunctionTok{aes}\NormalTok{(}\AttributeTok{x =}\NormalTok{ FPR, }\AttributeTok{y =}\NormalTok{ TPR)) }\SpecialCharTok{+}
  \FunctionTok{geom\_ribbon}\NormalTok{(}\FunctionTok{aes}\NormalTok{(}\AttributeTok{ymin =} \DecValTok{0}\NormalTok{, }\AttributeTok{ymax =}\NormalTok{ TPR), }\AttributeTok{fill =} \StringTok{"\#2980B9"}\NormalTok{, }\AttributeTok{alpha =} \FloatTok{0.15}\NormalTok{) }\SpecialCharTok{+}
  \FunctionTok{geom\_line}\NormalTok{(}\AttributeTok{color =} \StringTok{"\#2980B9"}\NormalTok{, }\AttributeTok{size =} \FloatTok{1.3}\NormalTok{) }\SpecialCharTok{+}
  \FunctionTok{geom\_abline}\NormalTok{(}\AttributeTok{slope =} \DecValTok{1}\NormalTok{, }\AttributeTok{intercept =} \DecValTok{0}\NormalTok{,}
              \AttributeTok{linetype =} \StringTok{"dashed"}\NormalTok{, }\AttributeTok{color =} \StringTok{"gray50"}\NormalTok{) }\SpecialCharTok{+}
  \FunctionTok{annotate}\NormalTok{(}\StringTok{"text"}\NormalTok{, }\AttributeTok{x =} \FloatTok{0.60}\NormalTok{, }\AttributeTok{y =} \FloatTok{0.12}\NormalTok{,}
           \AttributeTok{label =} \FunctionTok{sprintf}\NormalTok{(}\StringTok{"AUC = \%.3f"}\NormalTok{, auc\_r),}
           \AttributeTok{size =} \DecValTok{5}\NormalTok{, }\AttributeTok{color =} \StringTok{"\#2980B9"}\NormalTok{, }\AttributeTok{fontface =} \StringTok{"bold"}\NormalTok{) }\SpecialCharTok{+}
  \FunctionTok{labs}\NormalTok{(}\AttributeTok{title =} \StringTok{"Curva ROC"}\NormalTok{,}
       \AttributeTok{x =} \StringTok{"1 − Especificidade (FPR)"}\NormalTok{,}
       \AttributeTok{y =} \StringTok{"Sensibilidade (TPR)"}\NormalTok{) }\SpecialCharTok{+}
  \FunctionTok{coord\_equal}\NormalTok{() }\SpecialCharTok{+}
  \FunctionTok{theme\_minimal}\NormalTok{(}\AttributeTok{base\_size =} \DecValTok{12}\NormalTok{)}

\FunctionTok{grid.arrange}\NormalTok{(p\_mc, p\_roc, }\AttributeTok{ncol =} \DecValTok{2}\NormalTok{)}

\FunctionTok{cat}\NormalTok{(}\FunctionTok{sprintf}\NormalTok{(}\StringTok{"}\SpecialCharTok{\textbackslash{}n}\StringTok{AUC: \%.4f — \%s}\SpecialCharTok{\textbackslash{}n}\StringTok{"}\NormalTok{, auc\_r,}
\NormalTok{    dplyr}\SpecialCharTok{::}\FunctionTok{case\_when}\NormalTok{(}
\NormalTok{      auc\_r }\SpecialCharTok{\textgreater{}=} \FloatTok{0.90} \SpecialCharTok{\textasciitilde{}} \StringTok{"Excelente"}\NormalTok{,}
\NormalTok{      auc\_r }\SpecialCharTok{\textgreater{}=} \FloatTok{0.80} \SpecialCharTok{\textasciitilde{}} \StringTok{"Bom"}\NormalTok{,}
\NormalTok{      auc\_r }\SpecialCharTok{\textgreater{}=} \FloatTok{0.70} \SpecialCharTok{\textasciitilde{}} \StringTok{"Aceitável"}\NormalTok{,}
      \ConstantTok{TRUE}          \SpecialCharTok{\textasciitilde{}} \StringTok{"Fraco — revisar o modelo"}
\NormalTok{    )))}
\end{Highlighting}
\end{Shaded}

\subsection{7.11 Previsão para nova
observação}\label{previsuxe3o-para-nova-observauxe7uxe3o}

\begin{Shaded}
\begin{Highlighting}[]
\CommentTok{\# ── Defina aqui a nova observação (um valor por preditora) ──────────────────}
\NormalTok{valores\_novos\_r }\OtherTok{\textless{}{-}} \FunctionTok{c}\NormalTok{(}\DecValTok{45}\NormalTok{, }\DecValTok{55}\NormalTok{)   }\CommentTok{\# altere conforme sua análise}
\CommentTok{\# ────────────────────────────────────────────────────────────────────────────}

\NormalTok{nova\_obs\_r }\OtherTok{\textless{}{-}} \FunctionTok{setNames}\NormalTok{(}
  \FunctionTok{as.data.frame}\NormalTok{(}\FunctionTok{t}\NormalTok{(valores\_novos\_r)),}
\NormalTok{  VARIAVEIS\_PRED}
\NormalTok{)}

\NormalTok{prob\_nova\_r  }\OtherTok{\textless{}{-}} \FunctionTok{predict}\NormalTok{(modelo, }\AttributeTok{newdata =}\NormalTok{ nova\_obs\_r, }\AttributeTok{type =} \StringTok{"response"}\NormalTok{)}
\NormalTok{classe\_nova\_r }\OtherTok{\textless{}{-}} \FunctionTok{ifelse}\NormalTok{(prob\_nova\_r }\SpecialCharTok{\textgreater{}=}\NormalTok{ LIMIAR\_DECISAO,}
                        \StringTok{"Sucesso (1)"}\NormalTok{, }\StringTok{"Fracasso (0)"}\NormalTok{)}

\ControlFlowTok{for}\NormalTok{ (v }\ControlFlowTok{in}\NormalTok{ VARIAVEIS\_PRED)}
  \FunctionTok{cat}\NormalTok{(}\FunctionTok{sprintf}\NormalTok{(}\StringTok{"\%s = \%.1f}\SpecialCharTok{\textbackslash{}n}\StringTok{"}\NormalTok{, v, nova\_obs\_r[[v]]))}
\FunctionTok{cat}\NormalTok{(}\FunctionTok{sprintf}\NormalTok{(}\StringTok{"}\SpecialCharTok{\textbackslash{}n}\StringTok{Probabilidade estimada : \%.4f}\SpecialCharTok{\textbackslash{}n}\StringTok{"}\NormalTok{, prob\_nova\_r))}
\FunctionTok{cat}\NormalTok{(}\FunctionTok{sprintf}\NormalTok{(}\StringTok{"Classificação (τ = \%.2f): \%s}\SpecialCharTok{\textbackslash{}n}\StringTok{"}\NormalTok{, LIMIAR\_DECISAO, classe\_nova\_r))}
\end{Highlighting}
\end{Shaded}

\begin{center}\rule{0.5\linewidth}{0.5pt}\end{center}

\section{8. Implementação em
Python}\label{implementauxe7uxe3o-em-python}

\begin{quote}
\textbf{Como usar:} ajuste as \textbf{6 variáveis} do bloco
\texttt{\#\ CONFIGURAÇÃO} abaixo. O código usa \texttt{statsmodels} para
inferência completa e \texttt{scikit-learn} para métricas de
classificação.
\end{quote}

\subsection{8.1 Dependências}\label{dependuxeancias}

\begin{Shaded}
\begin{Highlighting}[]
\CommentTok{\# Instale com: pip install numpy pandas matplotlib seaborn scikit{-}learn statsmodels scipy}
\ImportTok{import}\NormalTok{ os, sys, warnings}
\ImportTok{import}\NormalTok{ numpy             }\ImportTok{as}\NormalTok{ np}
\ImportTok{import}\NormalTok{ pandas            }\ImportTok{as}\NormalTok{ pd}
\ImportTok{import}\NormalTok{ matplotlib.pyplot }\ImportTok{as}\NormalTok{ plt}
\ImportTok{import}\NormalTok{ matplotlib.gridspec }\ImportTok{as}\NormalTok{ gridspec}
\ImportTok{import}\NormalTok{ seaborn           }\ImportTok{as}\NormalTok{ sns}
\ImportTok{from}\NormalTok{ scipy  }\ImportTok{import}\NormalTok{ stats}
\ImportTok{from}\NormalTok{ sklearn.linear\_model    }\ImportTok{import}\NormalTok{ LogisticRegression}
\ImportTok{from}\NormalTok{ sklearn.model\_selection }\ImportTok{import}\NormalTok{ train\_test\_split}
\ImportTok{from}\NormalTok{ sklearn.metrics         }\ImportTok{import}\NormalTok{ (}
\NormalTok{    confusion\_matrix, classification\_report,}
\NormalTok{    roc\_curve, auc, ConfusionMatrixDisplay, accuracy\_score}
\NormalTok{)}
\ImportTok{import}\NormalTok{ statsmodels.api }\ImportTok{as}\NormalTok{ sm}

\NormalTok{warnings.filterwarnings(}\StringTok{"ignore"}\NormalTok{)}
\NormalTok{plt.rcParams.update(\{}\StringTok{"figure.dpi"}\NormalTok{: }\DecValTok{130}\NormalTok{, }\StringTok{"font.size"}\NormalTok{: }\DecValTok{11}\NormalTok{\})}
\end{Highlighting}
\end{Shaded}

\subsection{8.2 Configuração e importação dos
dados}\label{configurauxe7uxe3o-e-importauxe7uxe3o-dos-dados-1}

\begin{Shaded}
\begin{Highlighting}[]
\CommentTok{\# ════════════════════════════════════════════════════════════}
\CommentTok{\#  CONFIGURAÇÃO — ajuste apenas estas variáveis}
\CommentTok{\# ════════════════════════════════════════════════════════════}
\NormalTok{CAMINHO\_CSV        }\OperatorTok{=} \StringTok{"dados.csv"}
\NormalTok{VARIAVEL\_RESPOSTA  }\OperatorTok{=} \StringTok{"y"}
\NormalTok{VARIAVEIS\_PRED     }\OperatorTok{=}\NormalTok{ [}\StringTok{"x1"}\NormalTok{, }\StringTok{"x2"}\NormalTok{]}
\NormalTok{LIMIAR\_DECISAO     }\OperatorTok{=} \FloatTok{0.5}
\NormalTok{PROPORCAO\_TESTE    }\OperatorTok{=} \FloatTok{0.30}
\NormalTok{SEMENTE            }\OperatorTok{=} \DecValTok{42}
\CommentTok{\# ════════════════════════════════════════════════════════════}

\CommentTok{\# ── Dados simulados (se o .csv não existir) ─────────────────────────────────}
\ControlFlowTok{if} \KeywordTok{not}\NormalTok{ os.path.exists(CAMINHO\_CSV):}
    \BuiltInTok{print}\NormalTok{(}\SpecialStringTok{f"Arquivo não encontrado — gerando dados simulados."}\NormalTok{)}
\NormalTok{    rng  }\OperatorTok{=}\NormalTok{ np.random.default\_rng(SEMENTE)}
\NormalTok{    n\_s  }\OperatorTok{=} \DecValTok{350}
\NormalTok{    x1\_s }\OperatorTok{=}\NormalTok{ rng.uniform(}\DecValTok{18}\NormalTok{, }\DecValTok{65}\NormalTok{, n\_s).}\BuiltInTok{round}\NormalTok{(}\DecValTok{1}\NormalTok{)}
\NormalTok{    x2\_s }\OperatorTok{=}\NormalTok{ rng.normal(}\DecValTok{50}\NormalTok{, }\DecValTok{10}\NormalTok{, n\_s).}\BuiltInTok{round}\NormalTok{(}\DecValTok{1}\NormalTok{)}
\NormalTok{    z\_s  }\OperatorTok{=} \FloatTok{0.07} \OperatorTok{*}\NormalTok{ x1\_s }\OperatorTok{+} \FloatTok{0.06} \OperatorTok{*}\NormalTok{ x2\_s }\OperatorTok{{-}} \FloatTok{7.0}
\NormalTok{    p\_s  }\OperatorTok{=} \DecValTok{1} \OperatorTok{/}\NormalTok{ (}\DecValTok{1} \OperatorTok{+}\NormalTok{ np.exp(}\OperatorTok{{-}}\NormalTok{z\_s))}
\NormalTok{    y\_s  }\OperatorTok{=}\NormalTok{ rng.binomial(}\DecValTok{1}\NormalTok{, p\_s, n\_s)}
\NormalTok{    pd.DataFrame(\{}\StringTok{"y"}\NormalTok{: y\_s, }\StringTok{"x1"}\NormalTok{: x1\_s, }\StringTok{"x2"}\NormalTok{: x2\_s\}).to\_csv(CAMINHO\_CSV, index}\OperatorTok{=}\VariableTok{False}\NormalTok{)}
    \BuiltInTok{print}\NormalTok{(}\SpecialStringTok{f"\textquotesingle{}}\SpecialCharTok{\{}\NormalTok{CAMINHO\_CSV}\SpecialCharTok{\}}\SpecialStringTok{\textquotesingle{} criado com }\SpecialCharTok{\{}\NormalTok{n\_s}\SpecialCharTok{\}}\SpecialStringTok{ observações."}\NormalTok{)}

\CommentTok{\# ── Leitura ──────────────────────────────────────────────────────────────────}
\NormalTok{raw }\OperatorTok{=}\NormalTok{ pd.read\_csv(CAMINHO\_CSV)}
\BuiltInTok{print}\NormalTok{(}\SpecialStringTok{f"Dataset: }\SpecialCharTok{\{}\NormalTok{raw}\SpecialCharTok{.}\NormalTok{shape[}\DecValTok{0}\NormalTok{]}\SpecialCharTok{\}}\SpecialStringTok{ linhas × }\SpecialCharTok{\{}\NormalTok{raw}\SpecialCharTok{.}\NormalTok{shape[}\DecValTok{1}\NormalTok{]}\SpecialCharTok{\}}\SpecialStringTok{ colunas"}\NormalTok{)}
\BuiltInTok{print}\NormalTok{(}\SpecialStringTok{f"Colunas: }\SpecialCharTok{\{}\StringTok{\textquotesingle{}, \textquotesingle{}}\SpecialCharTok{.}\NormalTok{join(raw.columns)}\SpecialCharTok{\}}\SpecialStringTok{"}\NormalTok{)}

\CommentTok{\# Validação}
\NormalTok{cols\_faltando }\OperatorTok{=}\NormalTok{ [c }\ControlFlowTok{for}\NormalTok{ c }\KeywordTok{in}\NormalTok{ [VARIAVEL\_RESPOSTA] }\OperatorTok{+}\NormalTok{ VARIAVEIS\_PRED}
                 \ControlFlowTok{if}\NormalTok{ c }\KeywordTok{not} \KeywordTok{in}\NormalTok{ raw.columns]}
\ControlFlowTok{if}\NormalTok{ cols\_faltando:}
\NormalTok{    sys.exit(}\SpecialStringTok{f"Colunas ausentes: }\SpecialCharTok{\{}\StringTok{\textquotesingle{}, \textquotesingle{}}\SpecialCharTok{.}\NormalTok{join(cols\_faltando)}\SpecialCharTok{\}}\CharTok{\textbackslash{}n}\SpecialStringTok{"}
             \StringTok{"Ajuste VARIAVEL\_RESPOSTA e VARIAVEIS\_PRED."}\NormalTok{)}

\NormalTok{dados }\OperatorTok{=}\NormalTok{ raw[[VARIAVEL\_RESPOSTA] }\OperatorTok{+}\NormalTok{ VARIAVEIS\_PRED].copy()}
\NormalTok{dados.rename(columns}\OperatorTok{=}\NormalTok{\{VARIAVEL\_RESPOSTA: }\StringTok{"y"}\NormalTok{\}, inplace}\OperatorTok{=}\VariableTok{True}\NormalTok{)}
\NormalTok{dados[}\StringTok{"y"}\NormalTok{] }\OperatorTok{=}\NormalTok{ dados[}\StringTok{"y"}\NormalTok{].astype(}\BuiltInTok{int}\NormalTok{)}
\NormalTok{n\_antes }\OperatorTok{=} \BuiltInTok{len}\NormalTok{(dados)}
\NormalTok{dados.dropna(inplace}\OperatorTok{=}\VariableTok{True}\NormalTok{)}
\ControlFlowTok{if} \BuiltInTok{len}\NormalTok{(dados) }\OperatorTok{\textless{}}\NormalTok{ n\_antes:}
    \BuiltInTok{print}\NormalTok{(}\SpecialStringTok{f"Removidas }\SpecialCharTok{\{}\NormalTok{n\_antes }\OperatorTok{{-}} \BuiltInTok{len}\NormalTok{(dados)}\SpecialCharTok{\}}\SpecialStringTok{ linha(s) com NA."}\NormalTok{)}

\NormalTok{N1, N0, N }\OperatorTok{=}\NormalTok{ dados[}\StringTok{"y"}\NormalTok{].}\BuiltInTok{sum}\NormalTok{(), (dados[}\StringTok{"y"}\NormalTok{] }\OperatorTok{==} \DecValTok{0}\NormalTok{).}\BuiltInTok{sum}\NormalTok{(), }\BuiltInTok{len}\NormalTok{(dados)}
\BuiltInTok{print}\NormalTok{(}\SpecialStringTok{f"Classe 1: }\SpecialCharTok{\{}\NormalTok{N1}\SpecialCharTok{\}}\SpecialStringTok{ (}\SpecialCharTok{\{}\DecValTok{100}\OperatorTok{*}\NormalTok{N1}\OperatorTok{/}\NormalTok{N}\SpecialCharTok{:.1f\}}\SpecialStringTok{\%)   Classe 0: }\SpecialCharTok{\{}\NormalTok{N0}\SpecialCharTok{\}}\SpecialStringTok{ (}\SpecialCharTok{\{}\DecValTok{100}\OperatorTok{*}\NormalTok{N0}\OperatorTok{/}\NormalTok{N}\SpecialCharTok{:.1f\}}\SpecialStringTok{\%)"}\NormalTok{)}
\end{Highlighting}
\end{Shaded}

\subsection{8.3 Análise
exploratória}\label{anuxe1lise-exploratuxf3ria-1}

\begin{Shaded}
\begin{Highlighting}[]
\NormalTok{fig }\OperatorTok{=}\NormalTok{ plt.figure(figsize}\OperatorTok{=}\NormalTok{(}\DecValTok{13}\NormalTok{, }\DecValTok{4}\NormalTok{))}
\NormalTok{gs  }\OperatorTok{=}\NormalTok{ gridspec.GridSpec(}\DecValTok{1}\NormalTok{, }\DecValTok{1} \OperatorTok{+} \BuiltInTok{len}\NormalTok{(VARIAVEIS\_PRED), figure}\OperatorTok{=}\NormalTok{fig)}

\CommentTok{\# Dispersão}
\NormalTok{ax0   }\OperatorTok{=}\NormalTok{ fig.add\_subplot(gs[}\DecValTok{0}\NormalTok{, }\DecValTok{0}\NormalTok{])}
\NormalTok{cores }\OperatorTok{=}\NormalTok{ \{}\DecValTok{0}\NormalTok{: }\StringTok{"\#C0392B"}\NormalTok{, }\DecValTok{1}\NormalTok{: }\StringTok{"\#27AE60"}\NormalTok{\}}
\ControlFlowTok{for}\NormalTok{ cls, grp }\KeywordTok{in}\NormalTok{ dados.groupby(}\StringTok{"y"}\NormalTok{):}
\NormalTok{    ax0.scatter(grp[VARIAVEIS\_PRED[}\DecValTok{0}\NormalTok{]],}
\NormalTok{                grp[VARIAVEIS\_PRED[}\DecValTok{1}\NormalTok{]] }\ControlFlowTok{if} \BuiltInTok{len}\NormalTok{(VARIAVEIS\_PRED) }\OperatorTok{\textgreater{}} \DecValTok{1} \ControlFlowTok{else}\NormalTok{ grp[VARIAVEIS\_PRED[}\DecValTok{0}\NormalTok{]],}
\NormalTok{                c}\OperatorTok{=}\NormalTok{cores[cls],}
\NormalTok{                label}\OperatorTok{=}\StringTok{"Sucesso (1)"} \ControlFlowTok{if}\NormalTok{ cls }\OperatorTok{==} \DecValTok{1} \ControlFlowTok{else} \StringTok{"Fracasso (0)"}\NormalTok{,}
\NormalTok{                alpha}\OperatorTok{=}\FloatTok{0.65}\NormalTok{, edgecolors}\OperatorTok{=}\StringTok{"white"}\NormalTok{, s}\OperatorTok{=}\DecValTok{45}\NormalTok{)}
\NormalTok{ax0.set\_xlabel(VARIAVEIS\_PRED[}\DecValTok{0}\NormalTok{])}\OperatorTok{;}\NormalTok{ ax0.set\_title(}\StringTok{"Dispersão por Classe"}\NormalTok{)}
\NormalTok{ax0.set\_ylabel(VARIAVEIS\_PRED[}\DecValTok{1}\NormalTok{] }\ControlFlowTok{if} \BuiltInTok{len}\NormalTok{(VARIAVEIS\_PRED) }\OperatorTok{\textgreater{}} \DecValTok{1} \ControlFlowTok{else} \StringTok{""}\NormalTok{)}
\NormalTok{ax0.legend(fontsize}\OperatorTok{=}\DecValTok{9}\NormalTok{)}

\CommentTok{\# Boxplots}
\ControlFlowTok{for}\NormalTok{ k, var }\KeywordTok{in} \BuiltInTok{enumerate}\NormalTok{(VARIAVEIS\_PRED):}
\NormalTok{    ax }\OperatorTok{=}\NormalTok{ fig.add\_subplot(gs[}\DecValTok{0}\NormalTok{, k }\OperatorTok{+} \DecValTok{1}\NormalTok{])}
\NormalTok{    grupos }\OperatorTok{=}\NormalTok{ [dados.loc[dados[}\StringTok{"y"}\NormalTok{]}\OperatorTok{==}\DecValTok{0}\NormalTok{, var].values,}
\NormalTok{              dados.loc[dados[}\StringTok{"y"}\NormalTok{]}\OperatorTok{==}\DecValTok{1}\NormalTok{, var].values]}
\NormalTok{    bp }\OperatorTok{=}\NormalTok{ ax.boxplot(grupos, patch\_artist}\OperatorTok{=}\VariableTok{True}\NormalTok{, widths}\OperatorTok{=}\FloatTok{0.5}\NormalTok{,}
\NormalTok{                    medianprops}\OperatorTok{=}\BuiltInTok{dict}\NormalTok{(color}\OperatorTok{=}\StringTok{"black"}\NormalTok{, linewidth}\OperatorTok{=}\DecValTok{2}\NormalTok{))}
    \ControlFlowTok{for}\NormalTok{ b, cor }\KeywordTok{in} \BuiltInTok{zip}\NormalTok{(bp[}\StringTok{"boxes"}\NormalTok{], [}\StringTok{"\#C0392B"}\NormalTok{, }\StringTok{"\#27AE60"}\NormalTok{]):}
\NormalTok{        b.set\_facecolor(cor)}\OperatorTok{;}\NormalTok{ b.set\_alpha(}\FloatTok{0.75}\NormalTok{)}
\NormalTok{    ax.set\_xticklabels([}\StringTok{"Fracasso (0)"}\NormalTok{, }\StringTok{"Sucesso (1)"}\NormalTok{], fontsize}\OperatorTok{=}\DecValTok{9}\NormalTok{)}
\NormalTok{    ax.set\_title(}\SpecialStringTok{f"Distribuição de }\SpecialCharTok{\{}\NormalTok{var}\SpecialCharTok{\}}\SpecialStringTok{"}\NormalTok{)}\OperatorTok{;}\NormalTok{ ax.set\_ylabel(var)}

\NormalTok{plt.suptitle(}\StringTok{"Análise Exploratória"}\NormalTok{, fontweight}\OperatorTok{=}\StringTok{"bold"}\NormalTok{, y}\OperatorTok{=}\FloatTok{1.01}\NormalTok{)}
\NormalTok{plt.tight\_layout()}
\NormalTok{plt.savefig(}\StringTok{"py\_eda.png"}\NormalTok{, bbox\_inches}\OperatorTok{=}\StringTok{"tight"}\NormalTok{)}\OperatorTok{;}\NormalTok{ plt.show()}
\end{Highlighting}
\end{Shaded}

\subsection{8.4 Divisão treino/teste e ajuste do
modelo}\label{divisuxe3o-treinoteste-e-ajuste-do-modelo-1}

\begin{Shaded}
\begin{Highlighting}[]
\NormalTok{X }\OperatorTok{=}\NormalTok{ dados[VARIAVEIS\_PRED].values}
\NormalTok{y }\OperatorTok{=}\NormalTok{ dados[}\StringTok{"y"}\NormalTok{].values}

\NormalTok{X\_tr, X\_te, y\_tr, y\_te }\OperatorTok{=}\NormalTok{ train\_test\_split(}
\NormalTok{    X, y, test\_size}\OperatorTok{=}\NormalTok{PROPORCAO\_TESTE,}
\NormalTok{    random\_state}\OperatorTok{=}\NormalTok{SEMENTE, stratify}\OperatorTok{=}\NormalTok{y}
\NormalTok{)}
\BuiltInTok{print}\NormalTok{(}\SpecialStringTok{f"Treino: }\SpecialCharTok{\{}\BuiltInTok{len}\NormalTok{(y\_tr)}\SpecialCharTok{\}}\SpecialStringTok{ obs.  |  Teste: }\SpecialCharTok{\{}\BuiltInTok{len}\NormalTok{(y\_te)}\SpecialCharTok{\}}\SpecialStringTok{ obs."}\NormalTok{)}

\CommentTok{\# ── statsmodels: inferência estatística completa ────────────────────────────}
\NormalTok{X\_tr\_sm }\OperatorTok{=}\NormalTok{ sm.add\_constant(X\_tr)}
\NormalTok{X\_te\_sm }\OperatorTok{=}\NormalTok{ sm.add\_constant(X\_te)}
\NormalTok{modelo\_sm }\OperatorTok{=}\NormalTok{ sm.Logit(y\_tr, X\_tr\_sm).fit(disp}\OperatorTok{=}\VariableTok{False}\NormalTok{)}
\BuiltInTok{print}\NormalTok{(modelo\_sm.summary())}
\end{Highlighting}
\end{Shaded}

\subsection{8.5 Coeficientes e odds
ratios}\label{coeficientes-e-odds-ratios-1}

\begin{Shaded}
\begin{Highlighting}[]
\NormalTok{params  }\OperatorTok{=}\NormalTok{ modelo\_sm.params}
\NormalTok{ep      }\OperatorTok{=}\NormalTok{ modelo\_sm.bse}
\NormalTok{ci      }\OperatorTok{=}\NormalTok{ modelo\_sm.conf\_int()}

\NormalTok{nomes }\OperatorTok{=}\NormalTok{ [}\StringTok{"Intercepto"}\NormalTok{] }\OperatorTok{+}\NormalTok{ VARIAVEIS\_PRED}
\NormalTok{df\_coef }\OperatorTok{=}\NormalTok{ pd.DataFrame(\{}
    \StringTok{"Coeficiente"}\NormalTok{ : params.values,}
    \StringTok{"Erro Padrão"}\NormalTok{ : ep.values,}
    \StringTok{"OR"}\NormalTok{          : np.exp(params.values),}
    \StringTok{"IC 2,5\% (OR)"}\NormalTok{: np.exp(ci.iloc[:, }\DecValTok{0}\NormalTok{].values),}
    \StringTok{"IC 97,5\% (OR)"}\NormalTok{:np.exp(ci.iloc[:, }\DecValTok{1}\NormalTok{].values),}
    \StringTok{"Valor{-}p"}\NormalTok{     : modelo\_sm.pvalues.values}
\NormalTok{\}, index}\OperatorTok{=}\NormalTok{nomes)}

\BuiltInTok{print}\NormalTok{(df\_coef.}\BuiltInTok{round}\NormalTok{(}\DecValTok{4}\NormalTok{).to\_string())}
\end{Highlighting}
\end{Shaded}

\subsection{8.6 Pseudo-R² de McFadden}\label{pseudo-ruxb2-de-mcfadden-2}

\begin{Shaded}
\begin{Highlighting}[]
\NormalTok{L\_star\_py }\OperatorTok{=}\NormalTok{ modelo\_sm.llf}
\NormalTok{N1\_tr     }\OperatorTok{=} \BuiltInTok{int}\NormalTok{(y\_tr.}\BuiltInTok{sum}\NormalTok{())}
\NormalTok{N0\_tr     }\OperatorTok{=} \BuiltInTok{len}\NormalTok{(y\_tr) }\OperatorTok{{-}}\NormalTok{ N1\_tr}
\NormalTok{N\_tr      }\OperatorTok{=} \BuiltInTok{len}\NormalTok{(y\_tr)}
\NormalTok{L0\_py     }\OperatorTok{=}\NormalTok{ N1\_tr}\OperatorTok{*}\NormalTok{np.log(N1\_tr) }\OperatorTok{+}\NormalTok{ N0\_tr}\OperatorTok{*}\NormalTok{np.log(N0\_tr) }\OperatorTok{{-}}\NormalTok{ N\_tr}\OperatorTok{*}\NormalTok{np.log(N\_tr)}

\NormalTok{r2\_mcf\_py }\OperatorTok{=} \OperatorTok{{-}}\NormalTok{L\_star\_py }\OperatorTok{/}\NormalTok{ L0\_py}
\BuiltInTok{print}\NormalTok{(}\SpecialStringTok{f"L* = }\SpecialCharTok{\{}\NormalTok{L\_star\_py}\SpecialCharTok{:.4f\}}\SpecialStringTok{   L0 = }\SpecialCharTok{\{}\NormalTok{L0\_py}\SpecialCharTok{:.4f\}}\SpecialStringTok{"}\NormalTok{)}
\BuiltInTok{print}\NormalTok{(}\SpecialStringTok{f"Pseudo{-}R² de McFadden: }\SpecialCharTok{\{}\NormalTok{r2\_mcf\_py}\SpecialCharTok{:.4f\}}\SpecialStringTok{"}\NormalTok{)}
\BuiltInTok{print}\NormalTok{(}\SpecialStringTok{f"(statsmodels nativo):   }\SpecialCharTok{\{}\NormalTok{modelo\_sm}\SpecialCharTok{.}\NormalTok{prsquared}\SpecialCharTok{:.4f\}}\SpecialStringTok{"}\NormalTok{)}

\NormalTok{avaliacao }\OperatorTok{=}\NormalTok{ (}\StringTok{"Fraco"}       \ControlFlowTok{if}\NormalTok{ r2\_mcf\_py }\OperatorTok{\textless{}} \FloatTok{0.20} \ControlFlowTok{else}
             \StringTok{"Bom"}         \ControlFlowTok{if}\NormalTok{ r2\_mcf\_py }\OperatorTok{\textless{}} \FloatTok{0.40} \ControlFlowTok{else}
             \StringTok{"Muito bom"}\NormalTok{)}
\BuiltInTok{print}\NormalTok{(}\SpecialStringTok{f"Avaliação: }\SpecialCharTok{\{}\NormalTok{avaliacao}\SpecialCharTok{\}}\SpecialStringTok{"}\NormalTok{)}
\end{Highlighting}
\end{Shaded}

\subsection{8.7 Teste da razão de
verossimilhanças}\label{teste-da-razuxe3o-de-verossimilhanuxe7as-1}

\begin{Shaded}
\begin{Highlighting}[]
\NormalTok{modelo\_nulo\_py }\OperatorTok{=}\NormalTok{ sm.Logit(y\_tr, np.ones(N\_tr)).fit(disp}\OperatorTok{=}\VariableTok{False}\NormalTok{)}
\NormalTok{G\_py     }\OperatorTok{=} \DecValTok{2} \OperatorTok{*}\NormalTok{ (modelo\_sm.llf }\OperatorTok{{-}}\NormalTok{ modelo\_nulo\_py.llf)}
\NormalTok{gl\_py    }\OperatorTok{=} \BuiltInTok{len}\NormalTok{(VARIAVEIS\_PRED)}
\NormalTok{pval\_py  }\OperatorTok{=} \DecValTok{1} \OperatorTok{{-}}\NormalTok{ stats.chi2.cdf(G\_py, df}\OperatorTok{=}\NormalTok{gl\_py)}

\BuiltInTok{print}\NormalTok{(}\SpecialStringTok{f"H0: todos os coeficientes = 0"}\NormalTok{)}
\BuiltInTok{print}\NormalTok{(}\SpecialStringTok{f"Estatística G:      }\SpecialCharTok{\{}\NormalTok{G\_py}\SpecialCharTok{:.4f\}}\SpecialStringTok{"}\NormalTok{)}
\BuiltInTok{print}\NormalTok{(}\SpecialStringTok{f"Graus de liberdade: }\SpecialCharTok{\{}\NormalTok{gl\_py}\SpecialCharTok{\}}\SpecialStringTok{"}\NormalTok{)}
\BuiltInTok{print}\NormalTok{(}\SpecialStringTok{f"Valor{-}p:            }\SpecialCharTok{\{}\NormalTok{pval\_py}\SpecialCharTok{:.2e\}}\SpecialStringTok{"}\NormalTok{)}
\BuiltInTok{print}\NormalTok{(}\StringTok{"Decisão:"}\NormalTok{, }\StringTok{"Rejeitar H0."} \ControlFlowTok{if}\NormalTok{ pval\_py }\OperatorTok{\textless{}} \FloatTok{0.05} \ControlFlowTok{else} \StringTok{"Não rejeitar H0."}\NormalTok{)}
\end{Highlighting}
\end{Shaded}

\subsection{8.8 Teste de Wald}\label{teste-de-wald-1}

\begin{Shaded}
\begin{Highlighting}[]
\NormalTok{w\_py  }\OperatorTok{=}\NormalTok{ (params }\OperatorTok{/}\NormalTok{ ep) }\OperatorTok{**} \DecValTok{2}
\NormalTok{pw\_py }\OperatorTok{=} \DecValTok{1} \OperatorTok{{-}}\NormalTok{ stats.chi2.cdf(w\_py, df}\OperatorTok{=}\DecValTok{1}\NormalTok{)}

\NormalTok{df\_wald }\OperatorTok{=}\NormalTok{ pd.DataFrame(\{}
    \StringTok{"Coeficiente"}\NormalTok{  : params.values.}\BuiltInTok{round}\NormalTok{(}\DecValTok{4}\NormalTok{),}
    \StringTok{"Erro Padrão"}\NormalTok{  : ep.values.}\BuiltInTok{round}\NormalTok{(}\DecValTok{4}\NormalTok{),}
    \StringTok{"Estatística W"}\NormalTok{: w\_py.values.}\BuiltInTok{round}\NormalTok{(}\DecValTok{4}\NormalTok{),}
    \StringTok{"Valor{-}p"}\NormalTok{      : pw\_py.values.}\BuiltInTok{round}\NormalTok{(}\DecValTok{6}\NormalTok{),}
    \StringTok{"Decisão"}\NormalTok{      : [}\StringTok{"Rejeitar H0 (*)"} \ControlFlowTok{if}\NormalTok{ p }\OperatorTok{\textless{}} \FloatTok{0.05} \ControlFlowTok{else} \StringTok{"Não rejeitar H0"}
                      \ControlFlowTok{for}\NormalTok{ p }\KeywordTok{in}\NormalTok{ pw\_py]}
\NormalTok{\}, index}\OperatorTok{=}\NormalTok{nomes)}

\BuiltInTok{print}\NormalTok{(}\StringTok{"H0: a\_i = 0  |  H1: a\_i ≠ 0}\CharTok{\textbackslash{}n}\StringTok{"}\NormalTok{)}
\BuiltInTok{print}\NormalTok{(df\_wald.to\_string())}
\end{Highlighting}
\end{Shaded}

\subsection{8.9 Previsão e avaliação no conjunto de
teste}\label{previsuxe3o-e-avaliauxe7uxe3o-no-conjunto-de-teste-1}

\begin{Shaded}
\begin{Highlighting}[]
\CommentTok{\# scikit{-}learn para previsão e métricas de classificação}
\NormalTok{modelo\_sk }\OperatorTok{=}\NormalTok{ LogisticRegression(max\_iter}\OperatorTok{=}\DecValTok{1000}\NormalTok{, random\_state}\OperatorTok{=}\NormalTok{SEMENTE)}
\NormalTok{modelo\_sk.fit(X\_tr, y\_tr)}

\NormalTok{prob\_te   }\OperatorTok{=}\NormalTok{ modelo\_sk.predict\_proba(X\_te)[:, }\DecValTok{1}\NormalTok{]}
\NormalTok{y\_pred\_py }\OperatorTok{=}\NormalTok{ (prob\_te }\OperatorTok{\textgreater{}=}\NormalTok{ LIMIAR\_DECISAO).astype(}\BuiltInTok{int}\NormalTok{)}

\NormalTok{acc\_py  }\OperatorTok{=}\NormalTok{ accuracy\_score(y\_te, y\_pred\_py)}
\NormalTok{err\_py  }\OperatorTok{=} \DecValTok{1} \OperatorTok{{-}}\NormalTok{ acc\_py}
\BuiltInTok{print}\NormalTok{(}\SpecialStringTok{f"Acurácia:     }\SpecialCharTok{\{}\NormalTok{acc\_py}\SpecialCharTok{:.4f\}}\SpecialStringTok{  (}\SpecialCharTok{\{}\DecValTok{100}\OperatorTok{*}\NormalTok{acc\_py}\SpecialCharTok{:.1f\}}\SpecialStringTok{\%)"}\NormalTok{)}
\BuiltInTok{print}\NormalTok{(}\SpecialStringTok{f"Taxa de erro: }\SpecialCharTok{\{}\NormalTok{err\_py}\SpecialCharTok{:.4f\}}\SpecialStringTok{  (}\SpecialCharTok{\{}\DecValTok{100}\OperatorTok{*}\NormalTok{err\_py}\SpecialCharTok{:.1f\}}\SpecialStringTok{\%)}\CharTok{\textbackslash{}n}\SpecialStringTok{"}\NormalTok{)}
\BuiltInTok{print}\NormalTok{(classification\_report(y\_te, y\_pred\_py,}
\NormalTok{                             target\_names}\OperatorTok{=}\NormalTok{[}\StringTok{"Fracasso (0)"}\NormalTok{, }\StringTok{"Sucesso (1)"}\NormalTok{]))}
\end{Highlighting}
\end{Shaded}

\subsection{8.10 Visualizações: matriz de confusão e curva
ROC}\label{visualizauxe7uxf5es-matriz-de-confusuxe3o-e-curva-roc-1}

\begin{Shaded}
\begin{Highlighting}[]
\NormalTok{cm\_py }\OperatorTok{=}\NormalTok{ confusion\_matrix(y\_te, y\_pred\_py)}
\NormalTok{VN, FP, FN, VP }\OperatorTok{=}\NormalTok{ cm\_py.ravel()}
\NormalTok{sens }\OperatorTok{=}\NormalTok{ VP }\OperatorTok{/}\NormalTok{ (VP }\OperatorTok{+}\NormalTok{ FN) }\ControlFlowTok{if}\NormalTok{ (VP }\OperatorTok{+}\NormalTok{ FN) }\OperatorTok{\textgreater{}} \DecValTok{0} \ControlFlowTok{else} \DecValTok{0}
\NormalTok{spec }\OperatorTok{=}\NormalTok{ VN }\OperatorTok{/}\NormalTok{ (VN }\OperatorTok{+}\NormalTok{ FP) }\ControlFlowTok{if}\NormalTok{ (VN }\OperatorTok{+}\NormalTok{ FP) }\OperatorTok{\textgreater{}} \DecValTok{0} \ControlFlowTok{else} \DecValTok{0}
\NormalTok{prec }\OperatorTok{=}\NormalTok{ VP }\OperatorTok{/}\NormalTok{ (VP }\OperatorTok{+}\NormalTok{ FP) }\ControlFlowTok{if}\NormalTok{ (VP }\OperatorTok{+}\NormalTok{ FP) }\OperatorTok{\textgreater{}} \DecValTok{0} \ControlFlowTok{else} \DecValTok{0}
\NormalTok{f1   }\OperatorTok{=} \DecValTok{2}\OperatorTok{*}\NormalTok{prec}\OperatorTok{*}\NormalTok{sens }\OperatorTok{/}\NormalTok{ (prec}\OperatorTok{+}\NormalTok{sens) }\ControlFlowTok{if}\NormalTok{ (prec}\OperatorTok{+}\NormalTok{sens) }\OperatorTok{\textgreater{}} \DecValTok{0} \ControlFlowTok{else} \DecValTok{0}

\BuiltInTok{print}\NormalTok{(}\SpecialStringTok{f"Sensibilidade: }\SpecialCharTok{\{}\NormalTok{sens}\SpecialCharTok{:.4f\}}\SpecialStringTok{  Especificidade: }\SpecialCharTok{\{}\NormalTok{spec}\SpecialCharTok{:.4f\}}\SpecialStringTok{"}\NormalTok{)}
\BuiltInTok{print}\NormalTok{(}\SpecialStringTok{f"Precisão:      }\SpecialCharTok{\{}\NormalTok{prec}\SpecialCharTok{:.4f\}}\SpecialStringTok{  F1{-}Score:       }\SpecialCharTok{\{}\NormalTok{f1}\SpecialCharTok{:.4f\}}\SpecialStringTok{"}\NormalTok{)}

\NormalTok{fig, axes }\OperatorTok{=}\NormalTok{ plt.subplots(}\DecValTok{1}\NormalTok{, }\DecValTok{2}\NormalTok{, figsize}\OperatorTok{=}\NormalTok{(}\DecValTok{11}\NormalTok{, }\FloatTok{4.5}\NormalTok{))}

\CommentTok{\# Matriz de confusão}
\NormalTok{ConfusionMatrixDisplay(cm\_py,}
\NormalTok{    display\_labels}\OperatorTok{=}\NormalTok{[}\StringTok{"Fracasso (0)"}\NormalTok{, }\StringTok{"Sucesso (1)"}\NormalTok{]).plot(}
\NormalTok{    ax}\OperatorTok{=}\NormalTok{axes[}\DecValTok{0}\NormalTok{], colorbar}\OperatorTok{=}\VariableTok{False}\NormalTok{, cmap}\OperatorTok{=}\StringTok{"Blues"}\NormalTok{)}
\NormalTok{axes[}\DecValTok{0}\NormalTok{].set\_title(}\SpecialStringTok{f"Matriz de Confusão}\CharTok{\textbackslash{}n}\SpecialStringTok{Acurácia = }\SpecialCharTok{\{}\NormalTok{acc\_py}\SpecialCharTok{:.3f\}}\SpecialStringTok{"}\NormalTok{)}

\CommentTok{\# Curva ROC}
\NormalTok{fpr\_py, tpr\_py, \_ }\OperatorTok{=}\NormalTok{ roc\_curve(y\_te, prob\_te)}
\NormalTok{auc\_py            }\OperatorTok{=}\NormalTok{ auc(fpr\_py, tpr\_py)}
\NormalTok{axes[}\DecValTok{1}\NormalTok{].fill\_between(fpr\_py, tpr\_py, alpha}\OperatorTok{=}\FloatTok{0.15}\NormalTok{, color}\Oper
